\documentclass[11pt]{article}
\usepackage[a4paper,margin=21mm]{geometry}
\usepackage{fontspec}
\setmainfont{Latin Modern Roman}
\newfontfamily\CJKfont{Noto Serif CJK SC}
\XeTeXlinebreaklocale "zh"
\XeTeXlinebreakskip=0pt plus 1pt
\usepackage{amsmath,amssymb,amsthm,amscd,graphicx,color}
\usepackage{xurl}
\usepackage[unicode,hidelinks]{hyperref}
\allowdisplaybreaks
\setlength\emergencystretch{3em}
\numberwithin{equation}{section}

\usepackage{tikz-cd}
\usepackage{mathrsfs}

\numberwithin{equation}{section}
\newtheorem{Theorem}{Theorem}[section]
\newtheorem*{Theorem*}{Theorem}
\newtheorem{Lemma}[Theorem]{Lemma}
\newtheorem{Proposition}[Theorem]{Proposition}
\newtheorem{Conjecture}[Theorem]{Conjecture}
\newtheorem*{Conjecture*}{Conjecture}
{\theoremstyle{definition}
\newtheorem{Definition}[Theorem]{Definition}
\newtheorem{Example}[Theorem]{Example}
\newtheorem{Remark}[Theorem]{Remark}
}

\usetikzlibrary{matrix}
\usetikzlibrary{arrows}

\def\im{\operatorname{im}}
\def\loc{\operatorname{loc}}
\def\hol{\operatorname{hol}}
\def\dt{\operatorname{dt}}
\def\g{\operatorname{g}}
\def\ann{\operatorname{ann}}
\def\GL{\operatorname{GL}}
\def\SL{\operatorname{SL}}
\def\SO{\operatorname{SO}}
\def\O{\operatorname{O}}
\def\Ch{\operatorname{Ch}}
\def\td{\operatorname{Td}}
\def\Tot{\operatorname{Tot}}
\def\dom{\operatorname{dom}}
\def\range{\operatorname{range}}
\def\codom{\operatorname{codom}}
\def\t{\operatorname{t}}
\def\obj{\operatorname{obj}}
\def\mor{\operatorname{mor}}
\def\Gr{\operatorname{Gr}}
\def\res{\operatorname{res}}
\def\supp{\operatorname{supp}}
\def\Range{\operatorname{range}}
\def\Vect{\operatorname{Vect}}
\def\tind{\operatorname{t-ind}}
\def\ind{\operatorname{ind}}
\def\aind{\operatorname{a-ind}}
\def\id{\operatorname{id}}
\def\symb{\operatorname{Symb}}
\def\proj{\operatorname{proj}}
\def\hom{\operatorname{Hom}}
\def\germ{\operatorname{germ}}
\def\pdo{\operatorname{\Psi DO}}
\def\pt{\operatorname{pt}}
\def\Spin{\operatorname{Spin}}
\def\cl{\operatorname{cl}}
\def\obj{\operatorname{obj}}
\def\Diff{\operatorname{Diff}}
\def\mor{\operatorname{mor}}
\def\ab{\operatorname{\textbf{Ab}}}
\def\abs{\operatorname{\textbf{AbS}}}
\def\comp{\operatorname{\textbf{Comp}}}
\def\lcomp{\operatorname{\textbf{LocComp}}}
\def\vector{\operatorname{\textbf{Vect}}}
\def\Alg{\operatorname{\textbf{Alg}}}
\def\ring{\operatorname{\textbf{CRing}}}
\def\vol{\operatorname{vol}}
\def\Ell{\operatorname{Ell}}
\def\orb{\operatorname{orbit}}
\def\op{\operatorname{op}}
\def\hotimes{\operatorname{\hat{\otimes}}}
\def\coker{\operatorname{coker}}
\def\fred{\operatorname{Fred}}
\def\index{\operatorname{index}}
\def\av{\operatorname{Av}}
\def\Op{\operatorname{Op}}
\def\diam{\operatorname{diam}}
\def\gv{\operatorname{gv}}
\def\ch{\operatorname{ch}}
\def\PSL{\operatorname{PSL}}
\def\rank{\operatorname{rank}}
\def\Nd{\operatorname{End}}
\def\Cliff{\operatorname{\mathbb{C}liff}}
\def\Cl{\operatorname{Cl}}
\def\CCl{{\mathbb{C}\ell}}
\def\Ad{\operatorname{Ad}}
\def\Div{\operatorname{div}}
\def\End{\operatorname{End}}
\def\tr{\operatorname{tr}}
\def\Tr{\operatorname{Tr}}
\def\Fr{\operatorname{Fr}}
\def\Euc{\operatorname{Euc}}
\def\ad{\operatorname{ad}}
\def\Vol{\operatorname{Vol}}
\def\tGG{{\tilde{\mathcal{G}}}}
\def\SD{{\slashed D}}
\def\SB{{\slashed B}}
\def\ho{{\hat{\otimes}}}
\def\sgn{\operatorname{sgn}}
\def\ev{\operatorname{ev}}
\def\Aut{\operatorname{Aut}}
\def\Fl{\operatorname{Fl}}
\def\Top{\operatorname{Top}}
\def\triv{\operatorname{triv}}
\def\tgamma{\operatorname{\tilde{\gamma}}}

%--------mathcal
\def\EE{\operatorname{\mathcal{E}}}
\def\FF{{\mathcal{F}}}
\def\BB{{\mathcal{B}}}
\def\CC{{\mathcal{C}}}
\def\KK{{\mathcal{K}}}
\def\LL{{\mathcal{L}}}
\def\WW{{\mathcal{W}}}
\def\AA{{\mathcal{A}}}
\def\CC{{\mathcal{C}}}
\def\MM{{\mathcal{M}}}
\def\DD{{\mathcal{D}}}
\def\UU{{\mathcal{U}}}
\def\VV{{\mathcal{V}}}
\def\TT{{\mathcal{T}}}
\def\HH{{\mathcal{H}}}
\def\GG{{\mathcal{G}}}
\def\PP{{\mathcal{P}}}
\def\NN{{\mathcal{N}}}
\def\YY{{\mathcal{Y}}}
\def\XX{{\mathcal{X}}}
\def\OO{{\mathcal{O}}}

%-------mathscr
\def\ES{{\mathscr{E}}}
\def\FS{{\mathscr{F}}}
\def\BS{{\mathscr{B}}}
\def\CS{{\mathscr{C}}}
\def\KS{{\mathscr{K}}}
\def\LS{{\mathscr{L}}}
\def\WS{{\mathscr{W}}}
\def\AS{{\mathscr{A}}}
\def\CS{{\mathscr{C}}}
\def\MS{{\mathscr{M}}}
\def\DS{{\mathscr{D}}}
\def\US{{\mathscr{U}}}
\def\TS{{\mathscr{T}}}
\def\HS{{\mathscr{H}}}
\def\GS{{\mathscr{G}}}
\def\PS{{\mathscr{P}}}
\def\NS{{\mathscr{N}}}
\def\XS{{\mathscr{X}}}
\def\YS{{\mathscr{Y}}}
\def\SS{{\mathscr{S}}}

%-------mathbb
\def\EB{{\mathbb{E}}}
\def\CB{{\mathbb{C}}}
\def\RB{{\mathbb{R}}}
\def\ZB{{\mathbb{Z}}}
\def\NB{{\mathbb{N}}}
\def\TB{{\mathbb{T}}}
\def\HB{{\mathbb{H}}}

%------mathfrak
\def\AF{{\mathfrak{A}}}
\def\BF{{\mathfrak{B}}}
\def\EF{{\mathfrak{E}}}
\def\XF{{\mathfrak{X}}}
\def\gf{{\mathfrak{g}}}
\def\hf{{\mathfrak{h}}}
\def\pf{{\mathfrak{p}}}
\def\af{{\mathfrak{a}}}
\def\kf{{\mathfrak{k}}}
\def\glf{{\mathfrak{gl}}}
\def\vpf{{\mathfrak{vp}}}
\def\vf{{\mathfrak{v}}}

% Vertical space in tables
\def\tsep#1{%
 {\@tempdima=\ht\strutbox\advance\@tempdima #1
 \vrule height \@tempdima depth 0pt width 0pt\nobreak\hspace{0pt}%
 }%
}
\def\bsep#1{%
 {\@tempdima=\dp\strutbox \advance\@tempdima #1
 \nobreak\hspace{0pt}\vrule height 0pt depth \@tempdima width 0pt
 }%
}

\begin{document}
\begin{center}\Large Smooth and unique jet and germ parallel transport for singularly foliated bundles with a complete smooth partial connection and enough conservation laws\end{center}
\noindent\textbf{Stable ID:} 1410\quad\textbf{Author:} Lachlan Ewen MacDonald\par
\noindent\textbf{Work:} The Holonomy Groupoids of Singularly Foliated Bundles\par
\noindent\textbf{Source:} \url{https://sigma-journal.com/2021/043/}\par
\noindent\textbf{Version:} Official SIGMA 17 (2021) article 043, DOI 10.3842/SIGMA.2021.043; journal PDF and native TeX acquired 2026-09-30, exact bytes pinned by SHA256\par
\noindent\textbf{Rights:} CC BY 4.0. Authors retain copyright. \url{https://creativecommons.org/licenses/by/4.0/}. Reuse and typesetting adaptation; original language and mathematical content preserved.\par
\noindent\textbf{Difficulty (prior selection preserved):} {\CJKfont H2: The author constructs parallel transport on diffeological pseudo-bundles of invariant jets and germs, where ordinary manifold flow theory alone does not establish germ-level uniqueness. The proof uses a germ tangent-space description, bracket-preserving flow conjugation, parameter-dependent domain control and explicit plot smoothness. These obligations go beyond the regular-foliation or finite-dimensional connection argument.}\par
\medskip\noindent\textbf{Editorial boundary note (not author text):} Selected author-highlighted Theorem4.10: unique parallel transport for every finite jet order, inverse-limit infinity order and germs, with its smooth lifting map. All manifold/sheaf/flow, singular foliation, projectable jet prolongation, functional/quotient/tangent diffeology, section-germ pseudo-bundle and leafwise-path conventions are retained. Complete current germ-tangent calculation3.9, parameter-flow-domain lemma4.11, bracket/flow/invariance formulas, germ uniqueness and full three-neighborhood plot proof are supplied. Standard parameter-dependent manifold flows and foliation invariance under its own flows remain precise author-cited prerequisites. Later holonomy quotient/hierarchy/functoriality and unrelated sheaf-topology constructions are unselected; no technical component is separately counted. Native TikZ-CD diagrams remain editable with original arrows and labels, and original source numbering is preserved.\par
\medskip\hrule\medskip\noindent\textbf{Author text begins below.}\par
\setcounter{section}{1}
% Original source lines 243--243
In Section~\ref{sec4} we generalise the notion of a singular foliation to a \textit{singularly foliated bundle}. In~the same way that regularly foliated bundles in the sense of Kamber and Tondeur~\cite{folbund} are defined in terms of a partial connection on the total space of the bundle, our singularly foliated bundles are defined in terms of what we call \textit{singular partial connections}. Roughly speaking, a~singular partial connection $\ell$ in a fibre bundle over a singular foliation allows us to lift vector fields from the foliation of the base to fields on the total space. We~show that singularly foliated bundles generalise singular foliations and regularly foliated bundles, and are associated to all bundles that are equivariant under the action of a Lie groupoid. Associated to any singularly foliated bundle $\pi_{B}\colon B\rightarrow M$ with foliation $\FF$ of the base are pseudo-bundles $\pi_{B}^{k,\FF}\colon \Gamma_{k}(\pi_{B})^{\FF}\rightarrow M$ of germs (if~$k=\g$) or jets (if $k\leq\infty$) of sections which are invariant, to $k^{\rm th}$ order, under flows of~$\FF$. These $\FF$-invariant jets/germs generalise the jets/germs of ``distinguished sections'' considered in~\cite{mac3}, and may be thought of as conservation laws for $\FF$. We~prove the following.

\setcounter{section}{1}
% Original source lines 263--289
\section{Background and prerequisites}\label{sec2}

\subsection{Notational conventions}

All manifolds and fibre bundles are assumed to be smooth, Hausdorff, without boundary and connected, and all maps thereof assumed to be smooth. Given any manifold $M$, we use $\XF_{M}$ to denote the sheaf of smooth vector fields on $M$, and $C^{\infty}_{M}$ the sheaf of smooth, real-valued functions on $M$. If $\pi_{B}\colon B\rightarrow M$ is a fibre bundle, then we denote by $\Gamma_{\pi_{B}}$ the sheaf of sections of~$\pi_{B}$. Given a smooth map $f\colon M\rightarrow N$ of manifolds, and sheaves $\SS_{M}$ on $M$ and $\SS_{N}$ on $N$, we~denote by $f_{!}\SS_{M}$ the pushforward of $\SS_{M}$ and by $f^{!}\SS_{N}$ the pullback of $\SS_{N}$~\cite[p.~65]{hartshorne}. Given a (possibly time-dependent) vector field $X = X(t,-)$ on an open set $\OO$ in a manifold $M$, we use $\Fl^{X}$ to denote its \textit{flow}. That is, for $x\in\OO$ and $t_{0}\in\RB$, $t\mapsto\Fl^{X}_{t,t_{0}}(x)$ is the unique solution to the initial value problem
\begin{gather*}
\frac{\rm d}{{\rm d}s}\bigg|_{t}f(s;x) = X(t,f(t;x)),\qquad f(t_{0};x) = x
\end{gather*}
defined by the vector field $X$~\cite[p.~236]{leesmth}. For time-independent $X$, $\Fl^{X}_{t,t_{0}}$ depends only on $t-t_{0}$ and will be denoted simply by $\Fl^{X}_{t-t_{0}}$. Vector fields are assumed time-independent unless otherwise stated.

\subsection{Singular foliations}

We begin by recalling the standard sheaf-theoretic definition of a singular foliation.

\begin{Definition}\label{singdef}
Let $M$ be a smooth manifold. A~\textit{singular foliation} on $M$ is a subsheaf $\FF\subset \XF_{M}$ of $C^{\infty}_{M}$-modules on $M$ for which the following hold.
\begin{enumerate}\itemsep=0pt
\item $\FF$ is \textit{closed under Lie brackets} in the sense that for every open set $\OO$ of $M$, $\FF(\OO)$ is closed under the Lie bracket of vector fields.
\item $\FF$ is \textit{locally finitely generated} in the sense that for each $x\in M$, there exists an open neighbourhood $\OO_{x}$ of $x$ and a finite family $X_{1},\dots,X_{k}$ of elements of $\FF(\OO_{x})$ such that $\FF(\OO_{x})$ is the $C^{\infty}_{M}(\OO_{x})$-span of the $X_{i}$.
\end{enumerate}
\end{Definition}

One can alternatively describe a singular foliation of a manifold $M$ as a locally finitely generated submodule of the compactly supported vector fields on $M$ which is closed under Lie brackets, as in~\cite{iakovos2}. By~\cite[Remark 1.8]{GarZam} these definitions are equivalent. One of the most important facts regarding singular foliations is the Stefan--Sussmann integration theorem~\cite{stefan, sussmann}.

\begin{Theorem}%\label{stefansussmann}
Let $M$ be a manifold with a singular foliation $\FF$. Then $\FF$ integrates to give a~decom\-position of $M$ into smoothly immersed submanifolds called leaves.
\end{Theorem}

\setcounter{Theorem}{4}
% Original source lines 318--502
\subsection{Jet bundles and prolongation}

We recall in this subsection some well-known theory of jet bundles, drawn primarily from~\cite{varbi, saunders}. Although the reader is likely familiar with this theory already, we include the following outline both to introduce our rather unconventional notation (which we choose for consistency with the bundles of germs to be introduced in Section~\ref{pseudobundlesofgerms}) and to point out the structures that will be~most relevant in our constructions.

\begin{Definition}
Let $\pi_{B}\colon B\rightarrow M$ be a fibre bundle, with fibre $F$, and let $k\geq0$. We~say that two local sections $\sigma_{1}$ and $\sigma_{2}$ of $\pi_{B}$ defined in a neighbourhood of $x\in M$ have the same~$k$-\textit{jet at}~$x$ if $\sigma_{1}(x) =b= \sigma_{2}(x)$, and for any local coordinate neighbourhood $\OO\cong\OO_{M}\times\OO_{F}$ of $b$, one~has
\begin{gather*}
\frac{\partial^{|I|}\sigma^{i}_{1}}{\partial x^{I}}(x) = \frac{\partial^{|I|}\sigma^{i}_{2}}{\partial x^{I}}(x)
\end{gather*}
for all $i = 1,\dots,\dim(F)$, and all multi-indices $I$ with $|I|\leq k$. Having the same $k$-jet at a point~$x$ is an equivalence relation on the set of local sections defined about $x$, and we denote the $k$-jet equivalence class of any such local section by $[\sigma]^{k}_{x}$.
\end{Definition}

The $k$-jets of local sections fit into a fibre bundle in a natural way.

\begin{Definition}
Let $\pi_{B}\colon B\rightarrow M$ be a fibre bundle, and let $k\geq 0$. For each $x\in M$, denote the set of all $k$-jets of local sections defined near $x$ by $\Gamma_{k}(\pi_{B})_{x}$, and define
\begin{gather*}
\Gamma_{k}(\pi_{B}):=\bigsqcup_{x\in M}\Gamma_{k}(\pi_{B})_{x},
\end{gather*}
with $\pi_{B}^{k}\colon \Gamma_{k}(\pi_{B})\rightarrow M$ denoting the canonical projection. Define $\pi_{B}^{0,k}\colon \Gamma_{k}(\pi_{B})\rightarrow B$ by
\begin{gather*}
\pi_{B}^{0,k}\big(x,[\sigma]^{k}_{x}\big):=\sigma(x),
\end{gather*}
and observe that any choice of local coordinate trivialisation $(x^{i},f^{\alpha})$ on $\OO$ about a point $b\in B$ determines coordinates
\begin{gather*}%\label{coordsjets}
\big(x^{i},f^{\alpha},f^{\alpha}_{i},\dots,f^{\alpha}_{I}\big):=\bigg(x^{i},f^{\alpha},\frac{\partial f^{\alpha}}{\partial x^{i}},\dots,\frac{\partial^{|I|}f^{\alpha}}{\partial x^{I}}\bigg)
\end{gather*}
on the set $\big(\pi^{0,k}_{B}\big)^{-1}(\OO)$. These coordinates give the projection $\pi_{k}\colon \Gamma_{k}(\pi_{B})\rightarrow M$ the structure of~a~fibre bundle, called the $k^{\rm th}$ \textit{order jet bundle} of $\pi_{B}$.
\end{Definition}

The jet bundles of a fibre bundle $\pi_{B}\colon B\rightarrow M$ admit projections $\pi_{B}^{k,l}\colon \Gamma_{l}(\pi_{B})\rightarrow \Gamma_{k}(\pi_{B})$ defined by
\begin{gather*}
\pi_{B}^{k,l}\big([\sigma]^{l}_{x}\big):=[\sigma]^{k}_{x},\qquad [\sigma]^{l}_{x}\in \Gamma_{l}(\pi_{B})
\end{gather*}
for any $l\geq k$, and these projections form a projective system. The~projective limit of this system, denoted $\Gamma_{\infty}(\pi_{B})$, inherits a natural smooth structure as a projective limit of manifolds \cite[Chapter One, Section A]{varbi} (equivalently via the projective limit diffeology~\cite[Section 1.39]{diffeology}). The~space $\Gamma_{\infty}(\pi_{B})$ is usually identified with the set of $\infty$-jets of local sections of $\pi_{B}$, and admits a canonical projection $\pi_{B}^{\infty}\colon \Gamma_{\infty}(\pi_{B})\rightarrow M$ given by
\begin{gather*}
\pi_{B}^{\infty}\big([\sigma]^{\infty}_{x}\big):=x,\qquad [\sigma]^{\infty}_{x}\in \Gamma_{\infty}(\pi_{B}).
\end{gather*}
One therefore obtains a hierarchy of jet bundles
\begin{gather*}
\Gamma_{\infty}(\pi_{B})\rightarrow\cdots\xrightarrow{\pi_{B}^{k,k+1}} \Gamma_{k}(\pi_{B})\xrightarrow{\pi_{B}^{k-1,k}}\cdots\xrightarrow{\pi_{B}^{1,2}} \Gamma_{1}(\pi_{B})\xrightarrow{\pi_{B}^{0,1}}B\xrightarrow{\pi_{B}}M.
\end{gather*}

Since we will be concerned primarily with singular foliations, which arise from families of~vec\-tor fields, we will need to know about vector fields on jet bundles. A~particularly important class of vector fields on fibre bundles, in which we will be primarily interested, is those that are \textit{projectable} in the following sense.

\begin{Definition}%\label{projectable}
Let $\pi_{B}\colon B\rightarrow M$ be a fibre bundle. A~vector field $X$ on $B$ is said to be \textit{pro\-jec\-table} if there is a vector field $(\pi_{B})_{*}(X)$ on $M$ for which
\begin{gather}\label{projectableeqn}
{\rm d}\pi_{B}\circ X = (\pi_{B})_{*}(X)\circ \pi_{B}
\end{gather}
on all of $B$. We~denote by $\XF_{\proj}(B)$ the set of projectable vector fields on $B$.
\end{Definition}

Projectable vector fields on a bundle prolong in a natural way to vector fields on the associated jet bundles.

\begin{Definition}
Let $\pi_{B}\colon B\rightarrow M$ be a fibre bundle, and let $X\in\XF_{\proj}(B)$ be a (possibly time-dependent) projectable vector field. For $1\leq k<\infty$, the \textit{$k$-jet prolongation} of $X$ is the vector field $\pf^{k}(X)\in\XF(\Gamma_{k}(\pi_{B}))$ defined by
\begin{gather*}
\pf^{k}(X)\big(x,[\sigma]^{k}_{x}\big):=\frac{\rm d}{{\rm d}t}\bigg|_{t=0}\Big(\Fl^{(\pi_{B})_{*}X}_{t,0}(x),\big[\Fl^{X}_{t,0}\circ \sigma\circ\Fl^{(\pi_{B})_{*}X}_{0,t}\big]^{k}_{\Fl^{(\pi_{B})_{*}X}_{t,0}(x)}\Big)
\end{gather*}
for all $[\sigma]^{k}_{x}\in \Gamma_{k}(\pi_{B})$. We~denote by $\XF_{\proj}(\Gamma_{k}(\pi_{B}))$ the image of $\pf^{k}$.
\end{Definition}

It will be useful later to note here that if $X$ is a time-independent projectable vector field on~a~bundle $\pi_{B}$, it is given in coordinates by $X = a^{i}\frac{\partial}{\partial x^{i}} + b^{\alpha}\frac{\partial}{\partial f^{\alpha}}$, where $a^{i}$ are smooth functions depending only on the $x^{i}$ while the $b^{\alpha}$ depend on both the $x^{i}$ and the $f^{\alpha}$. Then for any $1\leq k<\infty$, the $k$-jet prolongation $\pf^{k}(X)$ of $X$ is given in coordinates over a point $(x^{i},f^{\alpha},f^{\alpha}_{i},\dots)\in \Gamma_{k}(\pi_{B})$ by the formula (cf.~\cite[Theorem~2.36]{olver})
\begin{gather}\label{prolongation}
\pf^{k}(X) = a^{i}D_{i}^{(k)}+\sum_{|I|=k}\!\bigg(D^{(k)}_{I}b^{\alpha}-\sum_{J\subset I}\big(D^{(k)}_{I\setminus J}a^{i}\big)f^{\alpha}_{Ji}\bigg)\frac{\partial}{\partial f^{\alpha}_{I}}+\!\sum_{|I|=0}^{k-1}D^{(k)}_{I}(b^{\alpha}-a^{i}f^{\alpha}_{i})\frac{\partial}{\partial f^{\alpha}_{I}},
\end{gather}
where the sum over $J\subset I$ is a sum over all \textit{strict} subsets of the multi-index $J$, and where we have absorbed the constants arising from the symmetry of mixed partial derivatives into our notation as in~\cite[equation~(1.15)]{varbi}. Here $D^{(k)}_{I} = D^{(k)}_{i_{1}}\cdots D^{(k)}_{i_{k}}$, where $D^{(k)}_{i}$ is the \textit{total derivative operator} defined by $D^{(k)}_{i} = \frac{\partial}{\partial x^{i}}+\sum_{|I|=0}^{k-1}f^{\alpha}_{Ii}\frac{\partial}{\partial f^{\alpha}_{I}}$.

The argument of Proposition~\ref{germprol} can be used to deduce the following fact, which justifies our choice of notation in denoting the image of $\pf^{k}$ in $\XF(\Gamma_{k}(\pi_{B}))$ by $\XF_{\proj}(\Gamma_{k}(\pi_{B}))$.

\begin{Proposition}
Let $\pi_{B}\colon B\rightarrow M$ be a fibre bundle, and let $X$ be a projectable vector field on~$B$. Then the $\pf^{k}(X)$ are a projectable family, in the sense that
\begin{gather*}
{\rm d}\pi_{B}^{l,k}\circ \pf^{k}(X) = \pf^{l}(X)\circ\pi_{B}^{l,k}
\end{gather*}
on $\Gamma_{k}(\pi_{B})$ for all $l\leq k$.
\end{Proposition}

Projectability of the $\pf^{k}(X)$ guarantees (cf.~\cite[equation~(1.11)]{varbi}) that pointwise they admit a~pro\-jective limit, defining a section $\pf^{\infty}(X)$ of the projective limit tangent bundle of $\Gamma_{\infty}(\pi_{B})$ (cf.~\cite[Chapter One, Section B]{varbi}) for which the following holds.

\begin{Proposition}\label{towerfield1}
Let $\pi_{B}\colon B\rightarrow M$ be a fibre bundle. Then for each $l\leq k\leq\infty$ there is a~bracket-preserving homomorphism $\big(\pi_{B}^{l,k}\big)_{*}\colon \XF_{\proj}(\Gamma_{k}(\pi_{B}))\rightarrow\XF_{\proj}(\Gamma_{l}(\pi_{B}))$ and the diagram
\begin{center}
\begin{tikzcd}[column sep = huge]
& \XF_{\proj}(\Gamma_{\infty}(\pi_{B})) \ar[d]
\\
& \vdots \ar[d,"\left(\pi_{B}^{1,2}\right)_{*}"]
\\
& \XF_{\proj}(\Gamma_{1}(\pi_{B})) \ar[d,"\left(\pi_{B}^{0,1}\right)_{*}"]
\\
\XF_{\proj}(B) \ar[uuur,"\pf^{\infty}"] \ar[ur,"\pf^{1}"] \ar[r,"\id"] & \XF_{\proj}(B)
\end{tikzcd}
\end{center}
commutes.
\end{Proposition}

\subsection{Diffeology}

We recall in this subsection some basic objects of study in diffeology that will be relevant for our constructions. The~most comprehensive reference on diffeology is the wonderful book~\cite{diffeology} by P. Iglesias-Zemmour.

\begin{Definition}\label{diffeology}
A function $\varphi\colon U\rightarrow \mathcal{X}$ from an open subset $U$ of some finite-dimensional Euclidean space to a set $\XX$ is called a \textit{parametrisation}. A~\textit{diffeology} on a set $\XX$ is a family D of~para\-metrisations satisfying the following axioms.
\begin{enumerate}\itemsep=0pt
\item The family D contains all constant parametrisations.
\item If $\varphi\colon U\rightarrow \XX$ is a parametrisation such that every point $u\in U$ has an open neighbourhood $V\subset U$ for which $\varphi|_{V}$ is an element of D, then $\varphi$ itself is an element of D.
\item For every element $\varphi\colon U\rightarrow \XX$ of D, every open set $V$ of any finite-dimensional Euclidean space, and for every smooth function $f\colon V\rightarrow U$, the composite $\varphi\circ f\colon V\rightarrow \XX$ is contained in D.
\end{enumerate}
A set with a diffeology is called a \textit{diffeological space}, and the elements of the diffeology are called its \textit{plots}. If $\XX$ and $\YY$ are two diffeological spaces, then a function $f\colon \XX\rightarrow \YY$ is said to be \textit{smooth} if for every plot $\varphi\colon U\rightarrow \XX$ of $\XX$, the composite $f\circ\varphi\colon U\rightarrow \YY$ is a plot of $\YY$. A~smooth bijection of diffeological spaces is said to be a \textit{diffeomorphism} if it is smooth with smooth inverse.
\end{Definition}

Every manifold is a diffeological space, with diffeology constituted by the set of all parametrisations that are smooth in the usual sense. Moreover a map between manifolds is smooth in the manifold sense if and only if it is smooth in the diffeological sense. Thus the category of~mani\-folds and smooth maps is a full and faithful subcategory of the category of diffeological spaces and smooth maps.

Diffeologies can be pushed forward and pulled back by functions of sets. This fact will be invoked frequently for our constructions.

\begin{Definition}
Let $\XX$ and $\YY$ be sets, and let $f\colon \XX\rightarrow \YY$ be a function.
\begin{enumerate}\itemsep=0pt
\item If $\XX$ has a diffeology, then the \textit{pushforward diffeology induced by $f$} is defined by declaring a parametrisation $\varphi\colon U\rightarrow \YY$ to be a plot if and only if every $u\in U$ has an open neighbourhood $V\subset U$ such that either $\varphi|_{V}$ is constant, or equal to the composite $f\circ\psi$ for some plot $\psi\colon V\rightarrow \XX$ of $\XX$. The~map $f$ is said to be a \textit{subduction} if it is surjective and if $\YY$ is equipped with the pushforward diffeology induced by $f$.
\item If $\YY$ has a diffeology, then the \textit{pullback diffeology induced by $f$} is defined by declaring a~parametrisation $\psi\colon U\rightarrow \XX$ to be plot if and only if the composite $f\circ\psi\colon U\rightarrow \YY$ is a~plot of~$\YY$. The~map $f$ is said to be an \textit{induction} if it is injective and if $\XX$ is equipped with the pullback diffeology from $\YY$.
\end{enumerate}
The following special cases are of particular importance. Let $\XX$ be a diffeological space.
\begin{enumerate}\itemsep=0pt
\item If $\sim$ is any equivalence relation on $\XX$, then the \textit{quotient diffeology} $\XX/{\sim}$ is the pushfoward diffeology arising from the quotient $\XX\rightarrow \XX/{\sim}$.
\item If $S$ is any subset of $\XX$, then the \textit{subspace diffeology} on $S$ is the pullback diffeology arising from the inclusion $S\hookrightarrow \XX$.
\item If $\YY$ is any other diffeological space, then the \textit{product diffeology} on $\XX\times \YY$ is the smallest diffeology for which the projections onto the factors are subductions.
\end{enumerate}
Quotients, subspaces and products will always be assumed to be equipped with the respective diffeologies defined above unless otherwise stated.
\end{Definition}

One of the features of the category of diffeological spaces is that the set of all morphisms between any two objects in the category is itself an object.

\begin{Definition}
Let $\XX$ and $\YY$ be diffeological spaces, and denote by $C^{\infty}(\XX,\YY)$ the set of~all smooth maps $f\colon \XX\rightarrow \YY$. The~\textit{functional diffeology} on $C^{\infty}(\XX,\YY)$ is defined by declaring a~para\-metrisation $\tilde{f}\colon U\rightarrow C^{\infty}(\XX,\YY)$ to be a plot if and only if the associated map $U\times\XX\ni(u,x)\mapsto\tilde{f}(u)(x)\rightarrow\YY$ is smooth.
\end{Definition}

The familiar notion of a fibre bundle over a manifold has a far reaching generalisation to diffeological spaces. It is the flexibility afforded by this generalisation that permits most of the constructions in this paper.

\begin{Definition}%\label{pb}
A \textit{diffeological pseudo-bundle} is a subduction $\pi_{\BB}\colon \BB\rightarrow \XX$ of diffeological spaces. Such a pseudo-bundle is in particular called a \textit{diffeological vector pseudo-bundle} if each fibre of $\pi_{\BB}$ is a vector space for which the vector space operations are smooth with respect to the subspace diffeology, and such that the zero section is smooth.
\end{Definition}

A diffeological pseudo-bundle need not have fibres that are all diffeomorphic, and of course need not be locally trivial in any sense (see for instance Example~\ref{germfol}). An important subclass of~diffeological pseudo-bundles are \textit{diffeological bundles}, which have mutually diffeomorphic fibres and which are locally trivial under pullbacks by plots. Diffeological bundles are distinguished by the behaviour of their \textit{structure groupoids}. Before we give the definition of this object, let us record what we mean by diffeological categories and groupoids.

\begin{Definition}
Let $\CC$ be a small category, with object set identified as a subset of the morphism set via the map which sends each object to its associated identity morphism. We~say that~$\CC$ is a \textit{diffeological category} if its set of morphisms is equipped with a diffeology for which the range, source, and composition are all smooth. If $\CC$ is in addition a groupoid, whose inversion map is smooth, we call $\CC$ a \textit{diffeological groupoid}.
\end{Definition}

Let now $\pi_{\BB}\colon \BB\rightarrow \XX$ be a smooth surjection of diffeological spaces. Denote by $\Aut(\pi_{\BB})$ the groupoid with object set $\XX$, and with morphisms from $x$ to $y$ constituted by the set $\Diff(\BB_{x},\BB_{y})$ of all diffeomorphisms from the fibre $\BB_{x}$ over $x$ to the fibre $\BB_{y}$ over $y$, with the obvious range, source, inversion and composition. The~groupoid $\Aut(\pi_{\BB})$ admits a smallest diffeology, called the \textit{functional diffeology}, under which the evaluation map $\ev\colon \Aut(\pi_{\BB})\times_{s,\pi_{\BB}}\BB\ni(f,b)\mapsto f(b)\in \BB$ is smooth, and under which $\Aut(\pi_{\BB})$ is a diffeological groupoid (see~\cite[Section~8.7]{diffeology} for details).

\begin{Definition}\label{structuregpd}
 Let $\pi_{\BB}\colon \BB\rightarrow \XX$ be a surjection of diffeological spaces. Equipped with the functional diffeology, we refer to $\Aut(\pi_{\BB})$ as the \textit{structure groupoid} of the surjection $\pi_{\BB}$. We~say that $\pi_{\BB}$ is a \textit{diffeological bundle} if the characteristic map $(r,s)\colon \Aut(\pi_{\BB})\rightarrow \XX\times \XX$ is a sub\-duction.
\end{Definition}

Diffeological bundles, unlike general diffeological pseudo-bundles, have a typical fibre to which all other fibres are diffeomorphic, and the pullback of a diffeological bundle along any plot is locally trivial~\cite[p.~240]{diffeology}.

By definition, singular foliations arise from certain families of sections of tangent bundles. To~use diffeology to study singular foliations, therefore, we need a notion of tangent bundle for a diffeological space. A~number of definitions have been proposed for this purpose, which, while coincident for manifolds, do not coincide for general diffeological spaces (see~\cite{cw} for a detailed discussion). The~point of view that we find useful here, as in~\cite{mac3}, is that of \textit{internal} tangent spaces and bundles. The~paper~\cite{cw} provides a categorical definition of internal tangent spaces based on work of Hector~\cite{hector}, which may be summarised as follows.

\begin{Definition}\label{tangentbundle}
Let $\XX$ be a diffeological space and let $x\in \XX$. Denote by $p_{x}$ the set of all \textit{plots centered at $x$}, that is, plots $\varphi\colon U\rightarrow \XX$ such that $0\in U$ and $\varphi(0) = x$. Denote by $T_{0}\dom(\varphi)$ the tangent space at zero of the domain of any such plot, and by $\vec{v}_{\varphi}$ the image of $\vec{v}\in T_{0}\dom(\varphi)$ in the direct sum $\bigoplus_{\varphi\in p_{x}}T_{0}\dom(\varphi)$. The~\textit{internal tangent space of $\XX$ at $x$} is the quotient space~$T_{x}\XX$ of the direct sum
\begin{gather*}
\bigoplus_{\varphi\in p_{x}} T_{0}\dom(\varphi)
\end{gather*}
by the subspace generated by all vectors of the form $f_{*}(\vec{v})_{\varphi} - \vec{v}_{\varphi\circ f}$, where $\varphi\in p_{x}$, $f\colon \dom(f)\rightarrow\dom(\varphi)$ is any smooth map of open Euclidean domains with $0\in\dom(f)$ and $f(0) = 0$, and with $\vec{v}\in T_{0}\dom(f)$. The~class of an element $\vec{v}_{\varphi}$, $\vec{v}\in T_{0}\dom(\varphi)$, will be denoted $\varphi_{*}(\vec{v})$. By~\cite[Proposition~3.3]{cw}, elements of $T_{x}\XX$ can always be written as linear combinations of $\varphi_{*}({\rm d}/{\rm d}t)$ for 1-plots $\varphi\colon (-\epsilon,\epsilon)\rightarrow\XX$, which we will usually write as ${\rm d}/{\rm d}t|_{0}\varphi_{t}$.
\end{Definition}

Let $\XX$ be a diffeological space, and consider the set $T\XX:=\bigsqcup_{x\in \XX}T_{x}\XX$. For any plot \mbox{$\varphi\colon U\rightarrow \XX$} and for any $u\in U$, denote by $\tau_{u}\colon v\mapsto v+u$ the $u$-translation map, so that $\tau_{u}^{-1}(U)$ contains $0$ and $\varphi\circ\tau_{u}\colon \tau_{u}^{-1}(U)\rightarrow \XX$ is a plot centered at $\varphi(u)$. Define ${\rm d}\varphi\colon TU\rightarrow T\XX$ by the formula
\begin{gather*}
{\rm d}\varphi(u,\vec{v}):=(\varphi(u),(\varphi\circ\tau_{u})_{*}(\vec{v})),\qquad (u,\vec{v})\in TU.
\end{gather*}
These maps were first considered by Hector~\cite{hector}. Then there exists a smallest diffeology on $T\XX$, called the \textit{dvs diffeology}~\cite{cw}, for which the natural projection $\pi_{T\XX}\colon T\XX\rightarrow \XX$ is a diffeological vector pseudo-bundle, and which contains the parametrisations ${\rm d}\varphi\colon TU\rightarrow T\XX$ as plots.

\begin{Definition}
Let $\XX$ be a diffeological space. The~diffeological vector pseudo-bundle $\pi_{T\XX}$: $T\XX\rightarrow \XX$ is called the \textit{internal tangent bundle of $\XX$}.
\end{Definition}

The internal tangent bundle is functorial under smooth maps of diffeological spaces.

\begin{Definition}\label{differential}
Let $f\colon \XX\rightarrow \YY$ be a smooth map of diffeological spaces. For any plot $\varphi$ of $\XX$ centered at $x$, define
\begin{gather*}
{\rm d}f(x,\varphi_{*}(\vec{v})):=(f(x),(f\circ\varphi)_{*}(\vec{v})),\qquad x\in \XX
\end{gather*}
and extend by linearity to a map ${\rm d}f\colon T\XX\rightarrow T\YY$. Then ${\rm d}f$ is a smooth map~\cite[Proposition~4.8]{cw} called the \textit{pushforward} or \textit{differential} of $f$.
\end{Definition}


% Original source lines 512--683
\section{Diffeological constructions}\label{sec3}

\subsection{Pseudo-bundles of germs}\label{pseudobundlesofgerms}

In~\cite{mac3}, we introduced ``bundles of germs'' of sections of certain fibre bundles. This construction can be generalised easily as follows. Let $\pi_{B}\colon B\rightarrow M$ be a smooth fibre bundle over a smooth manifold $M$, and let $\SS$ be a sheaf of smooth sections of $\pi_{B}$. Assume that $\SS$ is \textit{locally nonempty}, in the sense that for each $x\in M$, we can find an open neighbourhood $\mathcal{O}$ of $x$ such that $\SS(\mathcal{O})$ is nonempty. Define the \textit{total space} $\SS_{\loc}$ of the sheaf $\SS$ as the union
\begin{gather*}
\SS_{\loc}:=\bigcup_{\mathcal{O}}\SS(\mathcal{O})
\end{gather*}
over all open sets $\mathcal{O}$ in $M$ of elements of $\SS(\mathcal{O})$. Thus $\SS_{\loc}$ is the set of all locally defined sections of $\pi_{B}$ that belong to some $\SS(\mathcal{O})$. The~arguments of~\cite[Section~1.63]{diffeology} can be used to show that a diffeology may be defined on $\SS_{\loc}$ as follows.

\begin{Proposition}\label{functprop}
Let $M$ be a manifold, and let $\SS$ be a sheaf of smooth sections of some fibre bundle $B$ over $M$. Declare a parametrisation $\tilde{\sigma}\colon U\rightarrow \SS_{\loc}$ to have the property \textit{funct} if for all $u_{0}\in U$ and $x_{0}\in\dom(\tilde{\sigma}(u_{0}))$, there exists an open neighbourhood $V\subset U$ of $u_{0}$ and an open neighbourhood $\OO\subset \dom(\tilde{\sigma}(u_{0}))$ of $x_{0}$ for which $\OO\subset\dom(\tilde{\sigma}(u))$ for all $u\in V$ and for which the map $V\times\OO\ni(u,x)\mapsto\tilde{\sigma}(u)(x)\in B$ is smooth. Then the collection of parametrisations with the property \textit{funct} defines a diffeology on $\SS_{\loc}$.
\end{Proposition}

\begin{proof}
Observe first that if $\tilde{\sigma}\colon U\ni u\mapsto \sigma\in \SS(\mathcal{O}')$ is a constant plot, then we can simply take $\OO:=\OO'$ and $V = U$ to see that $\tilde{\sigma}$ has property \textit{funct}. Therefore axiom 1 of Definition~\ref{diffeology} is satisfied. To~see that axiom 2 is satisfied, suppose that $\tilde{\sigma}\colon U\rightarrow\SS_{\loc}$ is a parametrisation such that for each $u_{0}\in U$, there exists an open neighbourhood $V$ of $u_{0}$ such that $\tilde{\sigma}|_{V}$ has property \textit{funct}. Then by definition $\tilde{\sigma}$ must itself have property \textit{funct}, so that axiom 2 is satisfied. Finally, to prove that axiom 3 is satisfied, suppose that $\tilde{\sigma}\colon U\rightarrow\SS_{\loc}$ has property \textit{funct} and that $U'$ is some open Euclidean domain with $\varphi\colon U'\rightarrow U$ a smooth function. Let $u_{0}'\in U'$ and denote $u_{0}:=\varphi(u_{0}')\in U$. Let $V$ be an open neighbourhood of $u_{0}$ in $U$, and let $\OO$ be an open subset of $\dom(\tilde{\sigma}(u_{0}))$ for which $\OO\subset\dom(\tilde{\sigma}(u))$ for all $u\in V$ and for which $V\times\OO\ni(u,x)\mapsto\tilde{\sigma}(u)(x)\in B$ is smooth. Then $V':=\varphi^{-1}(V)$ is an open neighbourhood of $u_{0}'$ in $U'$, the open set $\OO$ satisfies $\OO\subset\dom(\tilde{\sigma}\circ\varphi(u))$ for all $u\in U$ and $V'\times\OO\ni(u',x)\mapsto\tilde{\sigma}(\varphi(u'))(x)\in B$ is smooth. Therefore axiom 3 is satisfied also.
\end{proof}

\begin{Definition}\label{funl}
Let $M$ be a manifold, and $\SS$ a sheaf of smooth sections of some fibre bundle over $M$. Then the diffeology on $\SS_{\loc}$ given in Proposition~\ref{functprop} is called the \textit{functional diffeology} on $\SS_{\loc}$.
\end{Definition}

Let us now consider the diffeological subspace
\begin{gather*}
\Gamma_{\loc}(\SS):=\{(x,\sigma)\colon x\in\dom(\sigma)\}
\end{gather*}
of the diffeological product $M\times \SS_{\loc}$. There is then clearly a surjective, smooth map $\pi_{\Gamma_{\loc}(\SS)}$: $\Gamma_{\loc}(\SS)\rightarrow M$ defined by
\begin{gather*}
\pi_{\Gamma_{\loc}(\SS)}(x,\sigma):=x,\qquad (x,\sigma)\in\Gamma_{\loc}(\SS),
\end{gather*}
which is moreover a subduction. Indeed, for any plot $\tilde{x}\colon U\rightarrow M$ and for any $u\in U$ we can always find an open neighbourhood $V$ of $u$ in $U$ such that $\tilde{x}(V)$ is contained in some open neighbour\-hood~$\mathcal{O}$ of $\tilde{x}(u)$ for which $\SS(\mathcal{O})$ contains some element $\sigma$. Now defining $\rho\colon V\rightarrow \Gamma_{\loc}(\SS)$ by~simply
\begin{gather*}
\rho(v):=(\tilde{x}(v),\sigma),\qquad v\in V
\end{gather*}
we have that $\pi_{\Gamma_{\loc}(\SS)}\circ\rho = \tilde{x}$, making $\pi_{\Gamma_{\loc}(\SS)}$ a subduction as claimed. The~fibre $\Gamma_{\loc}(\SS)_{x}$ over any $x\in M$ is the nonempty space consisting of sections $\sigma$ of $\SS$ defined on some open neighbourhood of $x$, equipped with the functional diffeology of Definition~\ref{funl}. The~subduc\-tion~$\pi_{\Gamma_{\loc}(\SS)}$ is the first step on the way to defining a genuinely useful object. Our next example shows why $\pi_{\Gamma_{\loc}(\SS)}$ is too large to be of much use in its own right.

\begin{Example}\label{ex1}
Let $\FF$ be a singular foliation of $M$, and denote the corresponding sheaf by the same symbol. Each fibre $\Gamma_{\loc}(\FF)_{x}$ is then \textit{almost} a Lie algebra. Indeed, if~$X\in\FF(\OO_{1})$ and \mbox{$Y\in\FF(\OO_{2})$} contain~$x$ in their domains of definition, then on the open neighbourhood $\OO:=\OO_{1}\cap \OO_{2}$ of~$x$, the Lie bracket $[X,Y]\in\FF(\OO)$ is defined. There is, however, nothing special about the choice $\OO:=\OO_{1}\cap\OO_{2}$, and indeed $[X,Y]$ also makes sense on any open neighbourhood~$\OO'$ of~$x$ within~$\OO$ and technically defines a distinct element $[X,Y]|_{\OO'}\in\FF(\OO')$. One encounters essentially the same problem when trying to define vector space operations in $\Gamma_{\loc}(\FF)_{x}$. To~rectify this sort of problem we work instead with a quotient of $\Gamma_{\loc}(\FF)$.
\end{Example}

Let us again return to a sheaf $\SS$ of sections of a bundle $\pi_{B}\colon B\rightarrow M$. Let us denote the germ at $x$ of any local section $\sigma$ of $\pi_{B}$ defined in an open neighbourhood of $x$ by $[\sigma]_{x}^{\g}$. We~define an equivalence relation~$\sim_{\g}$ on $\Gamma_{\loc}(\SS)$ by declaring $(x,\sigma)\sim_{\g}(y,\eta)$ if and only if $x = y$ and $[\sigma]^{\g}_{x} = [\eta]^{\g}_{x}$. We~denote by $\Gamma_{\g}(\SS)$ the diffeological quotient of $\Gamma(\SS)$ by the equivalence relation~$\sim_{\g}$, and denote by $\pi_{\SS}^{\g}\colon \Gamma_{\g}(\SS)\rightarrow M$ the obvious surjection
\begin{gather*}
\pi_{\SS}^{\g}\big(x,[\sigma]^{\g}_{x}\big):=x,\qquad \big(x,[\sigma]^{\g}_{x}\big)\in\Gamma_{\g}(\SS).
\end{gather*}
Since both the quotient map $q\colon \Gamma_{\loc}(\SS)\rightarrow\Gamma_{\g}(\SS)$ and the projection $\Gamma_{\loc}(\SS)\rightarrow M$ are subductions, so too is the projection $\pi_{\SS}^{\g}\colon \Gamma_{\g}(\SS)\rightarrow M$. Thus $\pi_{\SS}^{\g}\colon \Gamma_{\g}(\SS)\rightarrow M$ is a diffeological pseudo-bundle.

\begin{Definition}
Let $\SS$ be a sheaf of sections of a fibre bundle $\pi_{B}\colon B\rightarrow M$. Then the subduction $\pi_{\SS}^{\g}\colon \Gamma_{\g}(\SS)\rightarrow M$ is called the \textit{pseudo-bundle of germs of $\SS$}. If in particular $\SS$ is the sheaf of all sections of $\pi_{B}$, then we denote the pseudo-bundle of germs of $\SS$ by simply $\pi_{B}^{\g}\colon \Gamma_{\g}(\pi_{B})\rightarrow M$.
\end{Definition}

In fact the pseudo-bundle of germs of the full sheaf of sections of a fibre bundle $\pi_{B}\colon B\rightarrow M$ is a diffeological bundle, in the sense that all its fibres are isomorphic to the single diffeological space $C_{\g,0}^{\infty}\big(\RB^{\dim(M)};F\big)$ of germs at zero of smooth functions from $\RB^{\dim(M)}$ into the typical fibre~$F$ of~$B$. This can be seen, for instance, using an associated bundle construction in a similar fashion to~\cite[Remark 5.7]{mac3}~-- one need only replace the ``distinguished functions'' considered therein with coordinate maps on $M$. Thus it is entirely reasonable to refer to $\Gamma_{\g}(\pi_{B})$ as the \textit{bundle} of germs of sections of $\pi_{B}$. Let us now study its relationship with the jet bundles of $\pi_{B}$.

For each $k\leq\infty$ we have a canonical projection $\pi_{B}^{k,\g}\colon \Gamma_{\g}(\pi_{B})\rightarrow \Gamma_{k}(\pi_{B})$ onto the $k^{\rm th}$ order jet bundle of $\pi_{B}$ defined by
\begin{gather*}
\pi_{B}^{k,\g}\big(x,[\sigma]^{\g}_{x}\big):=\big(x,[\sigma]^{k}_{x}\big),\qquad \big(x,[\sigma]^{\g}_{x}\big)\in\Gamma_{\g}(\pi_{B}).
\end{gather*}
The arguments of~\cite[Proposition~5.14]{mac3} show that these projections are smooth, and are compatible with the jet projections $\pi_{B}^{l,k}\colon \Gamma_{k}(\pi_{B})\rightarrow \Gamma_{l}(\pi_{B})$ in the sense that $\pi_{B}^{l,\g} = \pi_{B}^{l,k}\circ\pi_{B}^{k,\g}$ for all $l\leq k$. We~therefore have a tower
\begin{center}
\begin{tikzcd}[row sep = large]
& & & \Gamma_{\g}(\pi_{B}) \ar[d,"\pi_{B}^{\infty,\g}"] & & &
\\
& & & \Gamma_{\infty}(\pi_{B}) \ar[dl,"\pi_{B}^{k+1,\infty}"] \ar[dr,"\pi_{B}^{k,\infty}"'] \ar[drrr,"\pi_{B}^{0,\infty}"'] & & &
\\
& \cdots\ar[r] & \Gamma_{k+1}(\pi_{B}) \ar[rr,"\pi_{B}^{k,k+1}"'] & & \Gamma_{k}(\pi_{B}) \ar[r] & \cdots \ar[r] & B
\end{tikzcd}
\end{center}
of diffeological bundles which extends the usual well-known tower of jet bundles of $\pi_{B}$.

Now there is not an easily identifiable Lie bracket on the space of vector fields (that is, sections of the internal tangent bundle) on $\Gamma_{\g}(\pi_{B})$, however there \textit{does} exist a certain diffeological subspace of vector fields which carries a natural Lie bracket. These vector fields are those that are contained in the image of a germinal prolongation operator from projectable vector fields on~$B$ to vector fields on $\Gamma_{\g}(\pi_{B})$.

\begin{Definition}
Let $\pi_{B}\colon B\rightarrow M$ be a fibre bundle. For (possibly time-dependent) $X\in\XF_{\proj}(B)$, the vector field $\pf^{\g}(X)$ on $\Gamma_{\g}(\pi_{B})$ defined by the formula
\begin{gather*}
\pf^{\g}(X)\big(x,[\sigma]^{\g}_{x}\big):=\frac{\rm d}{{\rm d}t}\bigg|_{t=0} \Big(\Fl^{(\pi_{B})_{*}X}_{t,0}(x),\big[\Fl^{X}_{t,0}\circ\sigma \circ\Fl^{(\pi_{B})_{*}(X)}_{0,t}\big]^{\g}_{\Fl^{(\pi_{B})_{*}(X)}_{t,0}(x)}\Big)
\end{gather*}
is called the \textit{germinal prolongation of $X$}. The~associated linear map $\pf^{\g}\colon \XF_{\proj}(B)\rightarrow \XF(\Gamma_{\g}(\pi_{B}))$ is called the \textit{germinal prolongation operator}, and its image is denoted $\XF_{\proj}(\Gamma_{\g}(\pi_{B}))$.
\end{Definition}

Our next result relates the germinal prolongation operator to the jet prolongations of projectable vector fields, and can be seen as a justification of the nomenclature ``germinal prolongation".

\begin{Proposition}\label{germprol}
Let $\pi_{B}\colon B\rightarrow M$ be a fibre bundle. Then for each $k\leq\infty$, we have
\begin{gather*}
{\rm d}\pi_{B}^{k,\g}\circ \pf^{\g}(X) = \pf^{k}(X)\circ\pi_{B}^{k,\g}
\end{gather*}
for all $($possibly time-dependent$)$ $X\in\XF_{\proj}(B)$. Consequently the tower of prolongations of~Pro\-po\-si\-tion~$\ref{towerfield1}$ completes to a tower
\begin{center}
\begin{tikzcd}[column sep = large, row sep = large]
& &\XF_{\proj}(\Gamma_{\g}(\pi_{B})) \ar[d, "(\pi_{B}^{\infty,\g})_{*}"]
\\
& & \XF_{\proj}(\Gamma_{\infty}(\pi_{B})) \ar[d]
\\
& & \vdots \ar[d,"\big(\pi_{B}^{0,1}\big)_{*}"]
\\
\XF_{\proj}(B) \ar[uuurr,"\pf^{\g}"] \ar[uurr,"\pf^{\infty}"] \ar[rr,"\id"] & & \XF_{\proj}(B)
\end{tikzcd}
\end{center}
\end{Proposition}

\begin{proof}
For any $k<\infty$ and $X\in\XF_{\proj}(B)$, we have
\begin{align*}
{\rm d}\pi_{B}^{k,\g}\circ\pf^{\g}(X)\big(x,[\sigma]^{\g}_{x}\big)
&= \frac{\rm d}{{\rm d}t}\bigg|_{t=0}\pi_{B}^{k,\g}\Big(\Fl^{(\pi_{B})_{*}X}_{t,0}(x),
\big[\Fl^{X}_{t,0}\circ\sigma\circ\Fl^{(\pi_{B})_{*}(X)}_{0,t}\big]^{\g}_{\Fl_{t,0}^{(\pi_{B})_{*}(X)}(x)}\Big)
\\
&= \frac{\rm d}{{\rm d}t}\bigg|_{t=0}\Big(\Fl^{(\pi_{B})_{*}X}_{t,0}(x), \big[\Fl^{X}_{t,0}\circ\sigma\circ\Fl^{(\pi_{B})_{*}(X)}_{0,t}\big]^{k}_{\Fl^{(\pi_{B})_{*}(X)}_{t,0}(x)}\Big)
\\
&= \pf^{k}(X)\big(\pi_{B}^{k,\g}\big(x,[\sigma]^{\g}_{x}\big)\big)
\end{align*}
for all $\big(x,[\sigma]^{\g}_{x}\big)\in\Gamma_{\g}(\pi_{B})$. For $k=\infty$ the result follows from the universal property of the projective limit of the $\pi_{B}^{k}$.
\end{proof}

The injectivity of the jet prolongation operators, which is apparent from equation~\eqref{prolongation} toge\-ther with Proposition~\ref{germprol}, implies that the germinal prolongation operator $\pf^{\g}\colon \XF_{\proj}(B)\rightarrow\XF(\Gamma_{\g}(\pi_{B}))$ of a bundle $\pi_{B}\colon B\rightarrow M$ is injective. Consequently, on $\XF_{\proj}(\Gamma_{\g}(\pi_{B}))$ we have a Lie bracket that is well-defined by the formula
\begin{gather*}
[\pf^{\g}(X),\pf^{\g}(Y)]:=\pf^{\g}([X,Y]),\qquad X,Y\in\XF_{\proj}(B),
\end{gather*}
with respect to which each $\big(\pi_{B}^{k,\g}\big)_{*}$ is a homomorphism of Lie algebras. While we will not be making use of this feature in this article, we remark that it distinguishes $\XF_{\proj}(\Gamma_{\g}(\pi_{B}))$ as a~rather special subspace of $\XF(\Gamma_{\g}(\pi_{B}))$, which, like the vector fields of many other diffeological spaces~\cite{cw}, does not carry a natural Lie bracket in general.

It is in examples arising from singular foliations that one sees the justification for the terminology ``pseudo-bundle'' in that the fibres of a pseudo-bundle of germs need not be isomorphic in general.

\begin{Example}\label{germfol}
Consider the foliation $\FF$ of $\RB$ generated by the vector field $X:=f\frac{\partial}{\partial x}$, where $f$ is any smooth function on $\RB$ such that $f(x) = 0$ for all $x\leq 0$ and such that $f(x)\neq 0$ for all $x>0$. Then for any $x_{0}<0$, the fibre $\Gamma_{\g}(\FF)_{x_{0}}$ consists of only a single point (every multiple of $X$ is equal to zero in a neighbourhood of $x_{0}$), while the fibre $\Gamma_{\g}(\FF)_{y_{0}}$ for any $y_{0}>0$ is the infinite-dimensional diffeological space consisting of all germs of smooth functions defined near the point $y_{0}$, and $\Gamma_{\g}(\FF)_{0}$ is the infinite-dimensional diffeological space consisting of germs of~smooth functions defined near zero that vanish at zero and below.
\end{Example}

\begin{Remark}
The pseudo-bundles of germs of singular foliations may be of particular interest~-- as we alluded to in Example~\ref{ex1}, they are canonically pseudo-bundles of diffeological Lie algebras. To~be more precise, let $\FF$ be a singular foliation of a manifold $M$. For each $x\in M$, let $\Gamma_{\loc}(\FF)_{x}$ be the diffeological subspace of $\Gamma_{\loc}(\FF)$ consisting of those local sections whose domains con\-tain~$x$, and let $q_{x}\colon \Gamma_{\loc}(\FF)_{x}\rightarrow\Gamma_{\g}(\FF)_{x}$ be the quotient map. Then it is routine to show that the formulae
\begin{gather*}
[q_{x}(X),q_{x}(Y)]:=q_{x}([X,Y]),\qquad
q_{x}(X)+q_{x}(Y) =q_{x}(X+Y),\qquad
\alpha q_{x}(X):=q_{x}(\alpha X)
\end{gather*}
for $X,Y\in\Gamma(\FF)_{x}$ and $\alpha\in\RB$ define a Lie algebra structure on $\Gamma_{\g}(\FF)_{x}$ which is smooth with respect to subspace diffeology from $\Gamma_{\g}(\FF)$. Thus $\pi_{\FF}^{\g}\colon \Gamma_{\g}(\FF)\rightarrow M$ is a diffeological vector pseudo-bundle of Lie algebras.
\end{Remark}

The final result of this subsection elucidates the nature of the tangent spaces to bundles of~germs. We~refer to Definition~\ref{tangentbundle} for notation.

\begin{Proposition}\label{tangentgerm}
Let $\pi_{B}\colon B\rightarrow M$ be a fibre bundle, with vertical tangent bundle $VB\rightarrow B$, and let $\big(x,[\sigma]^{\g}_{x}\big)\in\Gamma_{\g}(\pi_{B})$. Let $\Gamma_{\g}(\sigma^{*}VB)_{x}$ denote the vector space of germs at $x$ of sections of~$\sigma^{*}VB$. Then the map
\begin{gather}\label{tangerm}
\frac{\rm d}{{\rm d}t}\bigg|_{0}\big(\gamma(t),[\sigma_{t}]^{\g}_{\gamma(t)}\big) \mapsto \bigg(\frac{\rm d}{{\rm d}t}\bigg|_{0}\gamma(t),\bigg[\frac{\rm d}{{\rm d}t}\bigg|_{0}\sigma_{t}\bigg]_{x}\bigg),
\end{gather}
defines a linear isomorphism from $T_{\left(x,[\sigma]^{\g}_{x}\right)}\Gamma_{\g}(\pi_{B})$ to the vector space $T_{x}M\oplus\Gamma_{\g}(\sigma^{*}VB)_{x}$. In~particular, if ${\rm d}/{\rm d}t|_{0}\big(\gamma(t),[\sigma_{t}]^{\g}_{\gamma(t)}\big) = 0$, then there exists a neighbourhood $\OO$ of $x$ on which one has
\begin{gather*}
\frac{\rm d}{{\rm d}t}\bigg|_{0}\sigma_{t}(y) = 0
\end{gather*}
for all $y\in\OO$.
\end{Proposition}

\begin{proof}
First note that by definition of the diffeology on $\Gamma_{\g}(\pi_{B})$, any 1-plot $(-\epsilon,\epsilon)\rightarrow\Gamma_{\g}(\pi_{B})$ can, for sufficiently small $\epsilon$, be guaranteed to be of the form $t\mapsto \big(\gamma(t),[\sigma_{t}]^{\g}_{\gamma(t)}\big)$ for some plots $t\mapsto\sigma_{t}$ of $\Gamma_{\loc}(\pi_{B})$ and $t\mapsto\gamma(t)$ of $M$. Now observe that any sum of the form
\begin{gather*}
\frac{\rm d}{{\rm d}t}\bigg|_{0}\big(\gamma(t),[\sigma_{t}]^{\g}_{\gamma(t)}\big) +\frac{\rm d}{{\rm d}t}\bigg|_{0}\big(\tilde{\gamma}(t),[\tilde{\sigma}_{t}]^{\g}_{\tilde{\gamma}(t)}\big)
\end{gather*}
in $T_{\left(x,[\sigma]^{\g}_{x}\right)}\Gamma_{\g}(\pi_{B})$ can be represented by a single time derivative. By the first remark of this proof, to see this it suffices to work in $\Gamma_{\loc}(\pi_{B})$. Letting $\tilde{\OO}$ be a small neighbourhood of $\sigma(x)$ covering a small neighbourhood $\OO$ of $x$, let $\varphi\colon \tilde{\OO}\rightarrow\OO\times\RB^{m}$ denote the composite of coordinates on the fibre of $B$ with a local trivialisation of $B$ about $x$ such that $\varphi(\sigma(x)) = (x,0)$. Then the formula
\begin{gather*}%\label{kappa}
\psi(r,s):=\varphi^{-1}(\varphi\circ\sigma_{r}+\varphi\circ\tilde{\sigma}_{s})
\end{gather*}
defines a 2-parameter plot of $\Gamma_{\loc}(\pi_{B})$. Letting $\iota,\,\tilde{\iota}\colon \RB\rightarrow\RB^{2}$ denote the maps $t\mapsto(t,0)$ and $t\mapsto(0,t)$ respectively, we have $\sigma = \psi\circ\iota$ and $\tilde{\sigma} =\psi\circ\tilde{\iota}$, and therefore
\begin{gather*}
\frac{\rm d}{{\rm d}t}\bigg|_{0}\!\sigma_{t}+\frac{\rm d}{{\rm d}t}\bigg|_{0}\!\tilde{\sigma}_{t} =(\psi\circ\iota)_{*}\bigg(\frac{\rm d}{{\rm d}t}\bigg)+(\psi\circ\tilde{\iota})_{*}\bigg(\frac{\rm d}{{\rm d}t}\bigg)\! = \psi_{*}\bigg(\iota_{*}\bigg(\frac{\rm d}{{\rm d}t}\bigg)+\tilde{\iota}_{*}\bigg(\frac{\rm d}{{\rm d}t}\bigg)\!\bigg) = \psi_{*}\bigg(\frac{\rm d}{{\rm d}r},\frac{\rm d}{{\rm d}s}\bigg)
\end{gather*}
in $T_{\sigma}\Gamma_{\g}(\pi_{B})$. Now observe that if $h\colon \RB\rightarrow\RB^{2}$ denotes the map $t\mapsto(t,t)$, then, defining $\kappa_{t}:=\psi(t,t)$, we have
\begin{gather*}
\frac{\rm d}{{\rm d}t}\bigg|_{0}\kappa_{t} = (\psi\circ h)_{*}\bigg(\frac{\rm d}{{\rm d}t}\bigg) = \psi_{*}h_{*}\bigg(\frac{\rm d}{{\rm d}t}\bigg) = \psi_{*}\bigg(\frac{\rm d}{{\rm d}r},\frac{\rm d}{{\rm d}s}\bigg) = \frac{\rm d}{{\rm d}t}\bigg|_{0}\sigma_{t}+\frac{\rm d}{{\rm d}t}\bigg|_{0}\tilde{\sigma}_{t}.
\end{gather*}
Higher sums can be dealt with via a similar argument. Thus equation~\eqref{tangerm} does indeed suffice to define a map $T_{\left(x,[\sigma]^{\g}_{x}\right)}\Gamma_{\g}(\pi_{B})\rightarrow T_{x}M\oplus\Gamma_{\g}(\sigma^{*}VB)_{x}$, whose linearity is clear, and which admits the obvious inverse
\begin{gather*}
\bigg(\frac{\rm d}{{\rm d}t}\bigg|_{0}\gamma(t),\bigg[\frac{\rm d}{{\rm d}t}\bigg|_{0}\sigma_{t}\bigg]^{\g}_{x}\bigg)\mapsto \frac{\rm d}{{\rm d}t}\bigg|_{0}\big(\gamma(t),[\sigma]^{\g}_{\gamma(t)}\big).\tag*{\qed}
\end{gather*}
\renewcommand{\qed}{}
\end{proof}

\setcounter{subsection}{2}\setcounter{Theorem}{13}
% Original source lines 765--815
\subsection{The leafwise path category}

In~\cite{mac3}, we introduced a diffeological version of the Moore path category for any regular folia\-tion~$(M,\FF)$. The~objects of this category are simply points in $M$, while the morphisms are smooth, leafwise paths which have \textit{sitting instants} in that they are constant in small neighbourhoods of their endpoints. Composition of morphisms in this category is simply concatenation of paths. In~\cite{VilGar1}, the authors introduce an analogous diffeological space for \textit{singular} foliations, however concatenation of paths in this space no longer defines a category. In~this section, we~intro\-duce a hybrid of these two approaches~-- a diffeological space of integral curves of vector fields defining a singular foliation, for which concatenation of paths defines an associative and smooth multiplication.

We begin by recalling the definition of the path category of a diffeological space from~\cite{mac3} (cf.~\cite{ptfunctor}).

\begin{Definition}%\label{pathcat}
Let $\XX$ be a diffeological space. The~\textit{path category of $\XX$} is the diffeological sub\-space $\PP(\XX)$ of the diffeological product $C^{\infty}(\RB_{\geq0},\XX)\times\RB_{\geq0}$ consisting of pairs $(\gamma,d)$ for which there exist neighbourhoods of $0$ and of $[d,\infty)$ in $\RB_{\geq0}$ on which $\gamma$ is constant. Path categories are functorial~-- given any smooth map $f\colon \XX\rightarrow \YY$ of diffeological spaces, the formula
\begin{gather*}
\PP(f)(\gamma,d):=(f\circ\gamma,d)
\end{gather*}
defines a smooth functor $\PP(f)\colon \PP(\XX)\rightarrow \PP(\YY)$ of diffeological categories.
\end{Definition}

Given any diffeological space $\XX$, range and source maps $r$ and $s$ mapping $\PP(\XX)\rightarrow \XX$ are defined respectively by $(\gamma,d)\mapsto \gamma(d)$ and $(\gamma,d)\mapsto\gamma(0)$, and whenever $r(\gamma_{2},d_{2}) = s(\gamma_{1},d_{1})$, we~define the product $(\gamma_{1}\gamma_{2},d_{1}+d_{2})$ of $(\gamma_{1},d_{1})$ and $(\gamma_{2},d_{2})$ by the formula
\begin{gather*}
\gamma_{1}\gamma_{2}(t):=
\begin{cases}
\gamma_{2}(t),&\text{for}\quad 0\leq t\leq d_{2},
\\
\gamma_{1}(t-d_{2}),&\text{for}\quad d_{2}\leq t<\infty.
\end{cases}
\end{gather*}
This product, together with the range and source maps, are smooth, so that $\PP(\XX)$ is a dif\-feo\-lo\-gi\-cal category~\cite[Proposition~3.22]{mac3}. Moreover~\cite[Proposition~3.23]{mac3} there is a smooth involution $\iota\colon \PP(\XX)\ni(\gamma,d)\mapsto\big(\gamma^{-1},d\big)\rightarrow\PP(\XX)$ defined by the formula
\begin{gather*}
\gamma^{-1}(t):=
\begin{cases}
\gamma(d-t),&\text{if}\quad 0\leq t\leq d,
\\
\gamma(0),&\text{for}\quad t\geq d.
\end{cases}
\end{gather*}
Under favourable circumstances, which will be explicated in this section, the involution $\iota$ des\-cends to a genuine inversion on certain diffeological quotients of $\PP(\XX)$, giving such quotients the structures of diffeological groupoids.

Suppose in particular that $(M,\FF)$ is a singularly foliated manifold, and let $\Gamma_{\g}(\FF)$ be the pseudo-bundle of germs of $\FF$. By connectedness of $\RB_{\geq0}:=[0,\infty)$, any smooth map $\tilde{\gamma}\colon [0,\infty)\rightarrow\Gamma_{\g}(\FF)$ has the form
\begin{gather}\label{notation}
\tilde{\gamma}(t) = \big(\gamma(t),[X(t)]^{\g}_{\gamma(t)}\big),\qquad t\in[0,\infty),
\end{gather}
where $\gamma\colon [0,\infty)\rightarrow M$ is a smooth curve, and where $X\colon [0,\infty)\rightarrow\Gamma_{\loc}(\FF)$ is a smooth function. We~will implicitly use the notation of equation~\eqref{notation} in what follows.

\begin{Definition}\label{leafwisepaths}
Let $(M,\FF)$ be a singularly foliated manifold. We~define the \textit{leafwise} or \textit{$\FF$-path category} $\PP(\FF)$ to be the diffeological subspace of $\PP(\Gamma_{\g}(\FF))$ consisting of triples $\big(\gamma,[X]^{\g},d\big)$ for which $X\colon [0,\infty)\rightarrow\Gamma_{\loc}(\FF)$ satisfies
\begin{enumerate}\itemsep=0pt
\item[(1)] $\dom(X(t))$ is equal to a fixed open neighbourhood of $\gamma([0,\infty)) = \gamma([0,d])$ for all $t$,
\item[(2)] $X(t)(\gamma(t)) = \dot{\gamma}(t)$ for all $t\in[0,\infty)$, and
\item[(3)] $[X(0)]^{\g}_{\gamma(0)} = 0$ and $[X(t)]^{\g}_{\gamma(t)} = 0$ for all $t\geq d$.
\end{enumerate}
With range and source onto $M$ given by $r\big(\gamma,[X]^{\g},d\big):=\gamma(d)$ and $s\big(\gamma,[X]^{\g},d\big):=\gamma(0)$ respectively, by item 3, $\PP(\FF)$ inherits composition from $\PP(\Gamma_{\g}(\FF))$ to become a diffeological category with object space $M$.
\end{Definition}

Definition~\ref{leafwisepaths} is in practice the same as~\cite[Definition~3.1]{VilGar1} given by Garmendia--Villatoro. Note crucially that Definition~\ref{leafwisepaths} requires strictly more information than just the path in~$M$-re\-quiring in addition a time-dependent extension of the tangent field of $\gamma$ to an open neighbourhood of $\gamma$. This is so that flows of elements of $\PP(\FF)$ determine germs of diffeomorphisms defined in open neighbourhoods of their sources. Thus Definition~\ref{leafwisepaths} may be contrasted with the simpler~\cite[Definition~3.23]{mac3} for regular foliations, where such an extension is not explicitly required. In~the regular case, the tangent field along a leafwise path can always be \textit{canonically} extended to a tangent field in a (transverse) open neighbourhood of the path in any foliated chart. The~same is not true in the singular setting.


% Original source lines 841--1089
\section{Singularly foliated bundles and their holonomy groupoids}\label{sec4}

\subsection{Singularly foliated bundles}

Singular foliations are generalisations of regular foliations, and are naturally associated to Lie groupoids more generally. In~each of these special cases, one has a notion of fibre bundle which is compatible with the additional structure~-- in the case of a Lie groupoid action, the correct notion is that of an \textit{equivariant bundle}, while for a regular foliation the correct notion is that of a \textit{foliated bundle} in the sense of Kamber and Tondeur~\cite{folbund}. We~give in this section what appears to be the first definition of a fibre bundle compatible with a singular foliation, which simultaneously generalises equivariant and foliated bundles.

First, notice that projectable vector fields on a fibre bundle $\pi_{B}\colon B\rightarrow M$ over a manifold $M$ \textit{do not} generally form a sheaf of $C^{\infty}_{B}$-modules over $B$. Indeed, if $X$ is any projectable vector field and $f\in C^{\infty}(B)$ is any function which is non-constant along the fibres of $\pi_{B}$, then equation~\eqref{projectableeqn} will in general no longer hold for the vector field $fX$. Projectable vector fields are, however, closed under multiplication by functions of the form $f\circ\pi_{B}$, where $f\in C^{\infty}(M)$. Since $\pi_{B}$ is an open map, we can formulate the following definition, which will play a crucial role in our definition of singularly foliated bundle.

\begin{Definition}%\label{projectsheaf}
Let $\pi_{B}\colon B\rightarrow M$ be a fibre bundle. Denote by $C^{\infty}_{\proj,B}$ the subsheaf
\begin{gather*}
C^{\infty}_{\proj,B}(\OO):=\big\{f\circ\pi_{B}\in C^{\infty}_{B}(\OO)\colon f\in C^{\infty}_{M}(\pi_{B}(\OO))\big\}
\end{gather*}
of $C^{\infty}_{B}$, which we call the \textit{sheaf of projectable functions}. We~denote by $\XF_{\proj,B}$ the sheaf of $C^{\infty}_{\proj,B}$-modules
\begin{gather*}
\XF_{\proj,B}(\OO)\!:=\!\big\{X{\in}\XF_{B}(\OO)\colon \text{there is $(\pi_{B})_{*}X{\in}\XF_{M}(\pi_{B}(\OO))$ with ${\rm d}\pi_{B}\!\circ X\! = \! (\pi_{B})_{*}(X)\!\circ\!\pi_{B}$}\big\},
\end{gather*}
which we call the \textit{sheaf of projectable vector fields}.
\end{Definition}

The pushforward of projectable vector fields can now be characterised in the following sheaf-theoretic fashion. Recall that for a map $f\colon X\rightarrow Y$ of topological spaces and a sheaf $\SS$ on $X$, we use $f_{!}\SS$ to denote the pushforward of $\SS$ on $Y$~\cite[p.~65]{hartshorne}.

\begin{Proposition}
Let $\pi_{B}\colon B\rightarrow M$ be a fibre bundle. The~pushforward of projectable vector fields induces a morphism $(\pi_{B})_{*}\colon (\pi_{B})_{!}\XF_{\proj,B}\rightarrow\XF_{M}$ of sheaves of $C^{\infty}_{M}$-modules that preserves the Lie bracket.
\end{Proposition}

\begin{proof}
Notice first that we have a canonical isomorphism $C^{\infty}_{M}\cong (\pi_{B})_{!}C^{\infty}_{\proj,B}$ of sheaves of rings, obtained simply by sending $f\in C^{\infty}_{M}(\OO)$ to $f\circ\pi_{B}\in C^{\infty}_{\proj,B}(\pi_{B}^{-1}(\OO))$ for each open set $\OO$ in $M$. In~this way the $C^{\infty}_{\proj,B}$-module structure of $\XF_{\proj,B}$ indeed defines a $C^{\infty}_{M}$-module structure on $(\pi_{B})_{!}\XF_{\proj,B}$. The~pushforward $(\pi_{B})_{*}$ of $X\in\XF_{\proj,B}(\pi_{B}^{-1}(\OO))$ to $(\pi_{B})_{*}(X)\in\XF_{M}(\OO)$ then clearly preserves the associated $C^{\infty}_{M}(\OO)$-module structure for each open set $\OO$, and it is well-known~\cite[Lemma~3.10]{natopdiffgeom} that it also preserves the Lie bracket of vector fields.
\end{proof}

Singularly foliated bundles are now defined by \textit{singular partial connections}, which are particularly well-behaved partially-defined right-inverses of the pushforward morphism.

\begin{Definition}\label{singfolbund}
A \textit{singularly foliated bundle} is a triple $(\pi_{B},\FF,\ell)$, where $\pi_{B}\colon B\rightarrow M$ is a fibre bundle, $\FF$ is a singular foliation of $M$, and $\ell\colon \FF\rightarrow(\pi_{B})_{!}\XF_{\proj,B}$ is a morphism of sheaves of $C^{\infty}_{M}$-modules which preserves the Lie bracket, and for which the following hold.
\begin{enumerate}\itemsep=0pt
\item The morphism $\ell$ is a \textit{partial right-inverse} to $(\pi_{B})_{*}$ in the sense that
\begin{gather*}
(\pi_{B})_{*}\circ\ell = \id_{\FF},
\end{gather*}
on the sheaf $\FF$. In~particular this implies that $\ell$ is injective.
\item The morphism $\ell$ is \textit{complete} in the sense that for any open set $\OO$ in $M$, any $X\in\FF(\OO)$ and any $x\in\OO$, if $\Fl^{X}(x)$ is defined on an interval $I\subset\RB$ then so too is $\Fl^{\ell(X)}(b)$ for any $b\in B_{x}$.
\item The morphism $\ell$ is \textit{smooth} in the sense that the induced morphism
\begin{gather*}
\Gamma_{\loc}(\FF)\rightarrow \Gamma_{\loc}\big((\pi_{B})_{!}\XF_{\proj,B}\big)
\end{gather*}
of diffeological spaces is smooth with respect to the diffeology of Proposition~\ref{functprop}.
\end{enumerate}
We refer to such a morphism $\ell$ as a \textit{singular partial connection}.
\end{Definition}

We will usually denote a singularly foliated bundle $(\pi_{B},\FF,\ell)$ by simply $\pi_{B}$, with $\FF$ and $\ell$ assumed unless otherwise stated. Before discussing some examples, let us mention that completeness of $\ell$ does \textit{not} automatically follow from $\ell$ being a partial right inverse to $(\pi_{B})_{*}$. Indeed, it is easy to verify that for any open set $\OO\subset M$ and for any $X\in\FF(\OO)$, we have the relationship
\begin{gather*}
\Fl^{X}_{t}(x) = \pi_{B}\big(\Fl^{\ell(X)}_{t}(b)\big),\qquad x\in\OO,\ b\in B_{x}
\end{gather*}
between the flows of $\Fl^{X}(x)$ and $\Fl^{\ell(X)}(b)$ wherever they are defined. In~particular, that $\ell$ is a partial right-inverse to $(\pi_{B})_{*}$ implies that for any $b\in B_{x}$, the domain of $\Fl^{\ell(X)}(b)$ is contained in the domain of $\Fl^{X}(x)$. The~converse, however, does \textit{not} follow without completeness of $\ell$, as is easily seen by considering the standard example of $B =\RB^{2}$, $M = \RB$, and with $\ell\colon \XF(M)\rightarrow\XF_{\proj,B}(B)$ defined by
\begin{gather*}
\ell\bigg(f\frac{\partial}{\partial x}\bigg)(x,y):=f(x)\bigg(\frac{\partial}{\partial x}+y^{2}\frac{\partial}{\partial y}\bigg)
\end{gather*}
for $f\in C^{\infty}(M)$ and $(x,y)\in B$.

\begin{Example}[trivial bundles]\label{triv}
If $(M,\FF)$ is any singularly foliated manifold and $Q$ is any other manifold, then the trivial bundle $\pi\colon M\times Q\rightarrow M$ is canonically a singularly foliated bundle. Indeed, with respect to the decomposition $T(M\times Q)\cong TM\times TQ$, one has the trivial lift $\ell\colon \XF_{M}\rightarrow\pi_{!}\XF_{\proj,M\times Q}$ defined by the formula
\begin{gather*}
\ell(X):=(\pi^{*}(X),0)\in\XF_{M\times Q}\big(\pi^{-1}(\OO)\big) = (\pi_{!}\XF_{M\times Q})(\OO)
\end{gather*}
for all open sets $\OO$ in $M$ and $X\in\XF_{M}(\OO)$. Restricting $\ell$ to the subsheaf $\FF$ of $\XF_{M}$ one obtains a singular partial connection, with both completeness and smoothness being trivial.
\end{Example}

\begin{Example}[regularly foliated bundles]\label{regex}
Suppose that $\pi_{B}\colon B\rightarrow M$ is a \textit{regularly} foliated bundle, in the sense of Kamber--Tondeur~\cite[Definition~2.1]{folbund}~-- that is, there exists an involutive subbundle $T\FF_{B}\subset TB$ which is projected fibrewise-injectively to a subbundle $T\FF_{M}$ of $TM$. Involutivity of $T\FF_{B}$ implies that both $T\FF_{B}$ and $T\FF_{M}$ integrate to regular foliations of $B$ and $M$ respectively. Now if $\OO$ is any open subset of $M$ and $X\in\FF_{M}(\OO)$, then one obtains $\ell(X)\in\XF_{\proj,B}\big(\pi_{B}^{-1}(\OO)\big)$ whose value at a point $b\in B$ is the unique vector in $T_{b}\FF_{B}$ that is mapped by ${\rm d}\pi_{B}$ to $X(\pi_{B}(b))$. Clearly then the resulting morphism $\ell\colon \FF_{M}\rightarrow(\pi_{B})_{!}\XF_{\proj,B}$ is a singular partial connection in the sense of Definition~\ref{singfolbund}. Completeness and smoothness can both be seen by choosing foliated coordinates, in which $\ell$ is simply given by a trivial lift as in~Example~\ref{triv}.

Conversely, suppose that $\FF$ is a regular foliation of a manifold $M$, with leaf dimension $p$, and that $\pi_{B}\colon B\rightarrow M$ is a fibre bundle with a singular partial connection $\ell\colon \FF\rightarrow(\pi_{B})_{!}\XF_{\proj,B}$. In~a foliated chart $\OO\cong\RB^{p}\times\RB^{q}$ of $M$, wherein $\FF(\OO)$ is the $C^{\infty}_{M}(\OO)$-span of vector fields $\{e_{1},\dots,e_{p}\}$ that form the standard frame field of $\RB^{p}$, injectivity of $\ell$ implies that the vector fields $\{\ell(e_{1}),\dots,\ell(e_{p})\}$ span a $p$-dimensional subspace of $T_{b}B$ at each point $b\in\pi_{B}^{-1}(\OO)$ which intersects the vertical tangent space at $b$ only through zero. One thus obtains a smooth $p$-dimensional distribution $T\FF_{B}$ in $B$, and involutivity of $\FF$ together with the fact that $\ell$ preserves the Lie bracket implies that $T\FF_{B}$ is involutive. Thus $\pi_{B}\colon B\rightarrow M$ is a foliated bundle in~the sense of Kamber--Tondeur.
\end{Example}

\begin{Example}[equivariant bundles]\label{equiex}
Let $\GG$ be a Lie groupoid with unit space $M$. Denote the Lie algebroid of $\GG$ by $A:=\ker({\rm d}s)|_{M}$, with anchor map ${\rm d}r\colon A\rightarrow TM$, and let $\FF$ denote the associated sheaf of vector fields of the form ${\rm d}r\circ\sigma$, where $\sigma$ is an element of the sheaf of sections~$\AA$ of the Lie algebroid $A$. Let $\pi_{B}\colon B\rightarrow M$ be a fibre bundle.

If $\GG$ acts on $B$, then denote by $B\rtimes\GG$ the associated action groupoid~\cite[Example 2.2]{aac1}, with range and source denoted $r_{B}$ and $s_{B}$ respectively. Then the Lie algebroid $A_{B}:=\ker({\rm d}s_{B})|_{B}$ associated with the action groupoid $B\rtimes\GG$ is isomorphic to $\pi_{B}^{*}(A)$. For any open subset $\OO$ of~$M$, the formula
\begin{gather*}
\tilde{\ell}({\rm d}r\circ\sigma):={\rm d}r_{B}\circ\pi_{B}^{*}(\sigma),\qquad \sigma\in\AA(\OO)
\end{gather*}
then defines a singular partial connection. Completeness and smoothness are consequences of~the fact that $\tilde{\ell}$ is defined in terms of a smooth action of the groupoid $\GG$.
\end{Example}


\subsection{The holonomy groupoids of singularly foliated bundles}

One of the key objects in our construction of the holonomy groupoids of a singularly foliated bundle are pseudo-bundles of \textit{invariant} germs/jets. To~define these pseudo-bundles, we need \textit{vertical} prolongations of projectable vector fields (cf.~\cite[Definition~1.15]{varbi}).

\begin{Definition}
Let $\pi_{B}\colon B\rightarrow M$ be a fibre bundle, let $X\in\XF_{\proj}(B)$, and let $k$ denote any of~the symbols $1,\dots,\infty,\g$. The~vector field $\vpf^{k}(X)$ on $\Gamma_{k}(\pi_{B})$ defined by
\begin{gather*}
\vpf^{k}(X)\big(x,[\sigma]^{k}_{x}\big):=\frac{\rm d}{{\rm d}t}\bigg|_{t=0}\Big(x,\big[\Fl^{X}_{t}\circ\sigma\circ\Fl^{(\pi_{B})_{*}(X)}_{-t}\big]_{x}^{k}\Big)
\end{gather*}
is called the \textit{vertical $k$-prolongation} of $X$. On $\Gamma_{0}(\pi_{B}) = B$, we set $\vpf^{0}(X):=0$.
\end{Definition}

For a projectable vector field $X$ on a fibre bundle $\pi_{B}\colon B\rightarrow M$ and $k\geq1$, the vertical $k$-prolongation $\vpf^{k}(X)$ is the image of $\pf^{k}(X)$ under the canonical $\big(\pi_{B}^{k-1,k}\big)^{*}\big(V\pi_{B}^{k-1}\big)$-valued contact form $\theta^{(k)}$~\cite[Chapter~6.3]{saunders} on $\Gamma_{k}(\pi_{B})$. In~local coordinates $(x^{i},f^{\alpha},f^{\alpha}_{i},\dots)$, the components of~$\theta^{(k)}$ are given by
\begin{gather*}
\big(\theta^{(k)}\big)^{\alpha}_{I} = {\rm d}f^{\alpha}_{I} - f^{\alpha}_{Ii}{\rm d}x^{i}
\end{gather*}
for each $|I|<k$. Thus it is easily checked (cf.~equation~\eqref{prolongation}) that in coordinates the vertical prolongation of $X = a^{i}\frac{\partial}{\partial x^{i}}+b^{\alpha}\frac{\partial}{\partial f^{\alpha}}$ is given by
\begin{gather*}
\vpf^{k}(X) = \sum_{|I|=0}^{k-1}D_{I}^{(k)}(b^{\alpha}-a^{i}f^{\alpha}_{i})\frac{\partial}{\partial f^{\alpha}_{I}}.
\end{gather*}
The invariant pseudo-bundles of a singularly foliated bundle are now defined as follows.

\begin{Definition}\label{FFinv}
Let $\pi_{B}\colon B\rightarrow M$ be a singularly foliated bundle, and let $k$ denote any of the symbols $0,1,\dots,\infty,\g$. An element $\big(x,[\sigma]^{k}_{x}\big)$ of $\pi_{B}$ is said to be \textit{$\FF$-invariant to order $k$} if
\begin{gather*}
\vpf^{k}\big(\ell(X)\big)\big(x,[\sigma]^{k}_{x}\big) = 0
\end{gather*}
for all $X\in\FF$ defined in a neighbourhood $x$. We~say that $\pi_{B}$ \textit{admits enough conservation laws to order $k$} if at each $x\in M$ there is $\big(x,[\sigma]^{k}_{x}\big)$ which is $\FF$-invariant. This being the case, we call the corresponding sub-pseudo-bundle $\Gamma_{k}(\pi_{B})^{\FF}\subset\Gamma_{k}(\pi_{B})$ consisting of invariant germs/jets the \textit{$k^{\rm th}$ order $\FF$-invariant pseudo-bundle} of $\pi_{B}$. We~denote by $\pi_{B}^{k,\FF}$ the restriction of $\pi_{B}^{k}$ to $\Gamma_{k}(\pi_{B})^{\FF}$.
\end{Definition}

Note that if a singularly foliated bundle admits enough conservation laws to order $k>0$, then it admits enough conservation laws to any $l\leq k$. Moreover, by definition, \textit{every} singularly foliated bundle admits enough conservation laws to order 0.

Recall now~\cite[Definition~2.10]{mac3} that if $\pi_{B}\colon B\rightarrow M$ is a \textit{regularly} foliated bundle, a locally-defined section $\sigma$ of~$\pi_{B}$ is said to be \textit{distinguished} if, about any point in its image, there exist foliated coordinates $(x_{\alpha},y_{\alpha},f_{\alpha})$, with $x_{\alpha}$ and $y_{\alpha}$ denoting the leafwise and transverse coordinates respectively in the base, and with $f_{\alpha}$ denoting coordinates in the fibre, with respect to which $\sigma = \sigma(y_{\alpha})$ is independent of~the leafwise coordinates. We~denote by $\DS_{\g}(\pi_{B})$ the diffeological bundle of~germs of~the sheaf of~distinguished sections, and by $\DS_{k}(\pi_{B})$ the bundle of~jets of~distinguished sections. The~next proposition says that the $\FF$-invariant pseudo-bundles of~Defi\-ni\-tion~\ref{FFinv} generalise the bundles of~distinguished sections appearing in the regular case, and that therefore our constructions recover those of~\cite{mac3} in the regular case. In~particular, since distinguished functions of~this sort furnish the (local) degree 0 characteristic cohomology classes of~regular foliations~\cite[Example 1]{bryantgriffiths1}, we feel justified in referring to the $\FF$-invariant germs/jets of~a singularly foliated bundle as its \textit{conservation laws}.

\begin{Proposition}%\label{reg}
Let $\pi_{B}\colon B\rightarrow M$ be a regularly foliated bundle, and let $k$ denote any of the symbols $0,\dots,\infty,\g$. Then the diffeological subspace $\Gamma_{k}(\pi_{B})^{\FF}$ coincides with the space $\DS_{k}(\pi_{B})$ of classes of distinguished sections.
\end{Proposition}

\begin{proof}
In foliated coordinates $(x_{\alpha},y_{\alpha},f_{\alpha})$ for $B$, corresponding to leafwise, transverse and fibre coordinates respectively, any element $X\in\FF$ is given by some $C^{\infty}_{M}$-linear combination
\begin{gather*}
X = a^{i}\frac{\partial}{\partial x_{\alpha}^{i}},
\end{gather*}
while $\ell(X)$ (see Example~\ref{regex}) is given by
\begin{gather*}
\ell(X) = \pi_{B}^{*}(a^{i})\frac{\partial}{\partial x_{\alpha}^{i}}.
\end{gather*}
Thus, in our coordinates $(x_{\alpha},y_{\alpha},f_{\alpha})\in\RB^{p}\times\RB^{q}\times\RB^{k}$, we have simply
\begin{gather*}
\Fl^{\ell(X)}_{t}(x_{\alpha},y_{\alpha},f_{\alpha}) = \big(\Fl^{X}_{t}(x_{\alpha}),y_{\alpha},f_{\alpha}\big)
\end{gather*}
for small $t$. It follows immediately that for any smooth function $\sigma\colon \RB^{p}\times\RB^{q}\rightarrow\RB^{k}$, the curve
\begin{gather*}
t\mapsto\big(\Fl^{\ell(X)}_{t}\circ(\id\times\sigma)\circ\Fl^{X}_{-t}\big)(x_{\alpha},y_{\alpha}) = \big(x_{\alpha},y_{\alpha},\sigma\big(\Fl^{X}_{t}(x_{\alpha}),y_{\alpha}\big)\big)
\end{gather*}
is constant in $t$ \textit{for all $X\in\FF$} if and only if $\sigma$ is constant in the $x$ coordinate. Thus $\Gamma_{\g}(\pi_{B})^{\FF} = \DS_{\g}(\pi_{B})$. Since foliated coordinates can always be used to extend an invariant jet to an invariant local section, one has $\Gamma_{k}(\pi_{B})^{\FF} = \DS_{k}(\pi_{B})$ for $k\leq\infty$ also.
\end{proof}

Defining lifting maps and leafwise transport functors for a singularly foliated bundle is now a simple matter of putting our definitions together.

\begin{Theorem}\label{liftsmooth}
Let $\pi_{B}\colon B\rightarrow M$ be a singularly foliated bundle. Let $k$ denote any of the symbols $0,\dots,\infty,\g$, and suppose that $\pi_{B}$ admits enough conservation laws to order $k$. Then for each $\big(\gamma,[X]^{\g},d\big)\in\PP(\FF)$, $\gamma(0) = x$, and for each $\big(x,[\sigma]^{k}_{x}\big)\in\Gamma_{k}(\pi_{B})^{\FF}_{x}$, the map
\begin{gather}\label{ivpsol}
t\mapsto\big(\gamma(t),\big[\Fl^{\ell(X)}_{t,0}\circ\sigma\circ\Fl^{X}_{0,t}\big]^{k}_{\gamma(t)}\big)
\end{gather}
is the unique solution to the initial value problem
\begin{gather}\label{ivp}
\frac{\rm d}{{\rm d}s}\bigg|_{t}f\big(s;x,[\sigma]^{k}_{x}\big) = \pf^{k}\big(\ell(X(t))\big)\big(f\big(t;x,[\sigma]^{k}_{x}\big)\big),\qquad
f\big(0;x,[\sigma]^{k}_{x}\big) = \big(x,[\sigma]^{k}_{x}\big)
\end{gather}
in the diffeological space $\Gamma_{k}(\pi_{B})^{\FF}$ for which $\pi_{B}^{k,\FF}\circ f\big({-};x,[\sigma]^{k}_{x}\big) = \gamma(-)$. Moreover, the lifting map $L\big(\pi_{B}^{k,\FF}\big)\colon \PP(\FF)\times_{s,\pi_{B}^{k,\FF}}\Gamma_{k}(\pi_{B})^{\FF}\rightarrow \PP\big(\Gamma_{k}(\pi_{B})^{\FF}\big)$ defined by
\begin{gather}\label{Leqn}
L\big(\pi_{B}^{k,\FF}\big)\big(\gamma,[X]^{\g},d;x,[\sigma]^{k}_{x}\big)(t):=
\big(\gamma(t),\big[\Fl^{\ell(X)}_{t,0}\circ\sigma\circ\Fl^{X}_{0,t}\big]^{k}_{\gamma(t)}\big),\qquad t\in[0,\infty)
\end{gather}
is smooth.
\end{Theorem}

Note that the expression on the right hand side of equation~\eqref{Leqn} only makes sense by the completeness assumption on the singular partial connection. To~prove Theorem~\ref{liftsmooth}, we require the following lemma.

\begin{Lemma}\label{technicallemma}
Let $(M,\FF)$ be a singular foliation. Suppose that $U$ is an open set in Euclidean space and $X\colon \RB_{\geq0}\times U\rightarrow\Gamma_{\loc}(\FF)$ and $\gamma\colon \RB_{\geq0}\times U\rightarrow M$ are smooth maps for which
\begin{enumerate}\itemsep=0pt
\item[$(1)$] $\dom(X(t,u)) = \dom(X(s,u))$ for all $s,t\in\RB_{\geq0}$ and $u\in U$, and for which
\item[$(2)$] $t\mapsto\gamma(t,u)$ is an integral curve of $X(t,u)$ for all $(t,u)\in\RB_{\geq0}\times U$. That is,
\begin{gather*}
\frac{\rm d}{{\rm d}s}\bigg|_{s=t}\gamma(s,u) = X(t,u)(\gamma(t,u))
\end{gather*}
for all $(t,u)\in\RB_{\geq0}\times U$.
\end{enumerate}
Then for each $t_{0}\in\RB_{\geq0}$ and $u_{0}\in U$, there exist $\epsilon>0$ and open neighbourhoods $\UU\ni u_{0}$ in $U$ and $\OO\ni x_{0}:=\gamma(0,u_{0})$ in $M$, such that $\OO\subset\dom\big(\Fl^{X(u)}_{t,0}\big)$ for all $(t,u)\in[0,t_{0}+\epsilon)\times\UU$, and for which the map
\begin{gather*}
[0,t_{0}+\epsilon)\times\UU\times\OO\ni(t,u,x)\mapsto\Fl^{X(u)}_{t,0}(x)\in M
\end{gather*}
is smooth.
\end{Lemma}

\begin{proof}
By the definition of the diffeology on $\Gamma_{\loc}(\FF)$ and hypothesis 1, we can find open sets $\UU'\ni u_{0}$ and $\OO'\ni x_{0}$ such that $\OO'\subset\dom(X(t,u))$ for all $(t,u)\in\RB_{\geq0}\times\UU'$, and on which
\begin{gather*}
\RB_{\geq0}\times \UU'\times\OO'\ni(t,u,x)\mapsto X(t,u)(x)\in TM
\end{gather*}
is smooth. By hypothesis 2 we can assume that $\OO'$ contains $\gamma(\RB_{\geq0}\times\{u_{0}\})$. Now $X$ may be regarded as a time-dependent vector field on the manifold $\UU'\times\OO'$, and then by standard theory~\cite[Theorem~9.48]{leesmth}, there exists some maximal open neighbourhood $\NN\subset\UU'\times\OO'$ of~$(u_{0},x_{0})$ on which $t\mapsto\Fl^{X(u)}_{t,0}(x)$ is defined for all $u\in\UU$, $x\in\OO$ and for small $t$. Since by hypothesis $t\mapsto\Fl^{X(u_{0})}_{t,0}(x_{0}) = \gamma(t,u_{0})$ is defined for all $t\in\RB_{\geq0}$ we can always choose $\epsilon>0$, $\UU\ni u_{0}$ and $\OO\ni x_{0}$ small enough that $(t,u,x)\mapsto\Fl^{X(u)}_{t,0}(x)$ is well-defined and smooth on $[0,t_{0}+\epsilon)\times\UU\times\OO$ as claimed.
\end{proof}

\begin{proof}[Proof of Theorem~\ref{liftsmooth}]
That equation~\eqref{ivpsol} defines a curve in $\Gamma_{k}(\pi_{B})^{\FF}$ follows from the inva\-riance of $\FF$ under its own (possibly time-dependent~\cite[Lemma~3.3]{VilGar1}) flows~\cite[Proposition~1.6]{iakovos2}, the $\FF$-invariance of $\sigma$, and that $\ell$ is bracket preserving. More precisely, given $t\in[0,\infty)$, an open neighbourhood $\OO$ of $\gamma(t)$ and $Y\in\FF(\OO)$, since $\ell$ is bracket preserving we have
\begin{gather}\label{bracketpres}
\ell([X(t),Y])\big(\Fl^{\ell(X)}_{t,0}(b)\big) = [\ell(X(t)),\ell(Y)]\big(\Fl^{\ell(X)}_{t,0}(b)\big)
\end{gather}
for all $b\in\Fl^{\ell(X)}_{0,t}(\OO)$. Now, for any such $b$, with $\pi_{B}(b) = x'$, the right hand side of equation~\eqref{bracketpres} is the lift by $\ell$ of the tangent to the curve $r\mapsto\big(\Fl^{X}_{0,t+r}\big)_{*}\big(Y\big(\Fl_{t+r,0}^{X}(x')\big)\big)$ at $r = 0$,
while the left hand side is the tangent to the curve $r\mapsto \big(\Fl^{\ell(X)}_{0,t+r}\big)_{*}\big(\ell(Y)\big(\Fl^{\ell(X)}_{t+r,0}(b)\big)\big)$ at $r=0$. By uniqueness of flows therefore, we have
\begin{gather}\label{ellhom}
\big(\Fl^{\ell(X)}_{0,t}\big)_{*}(\ell(Y)) = \ell\big(\big(\Fl^{X}_{0,t}\big)_{*}(Y)\big).
\end{gather}
For notational simplicity denote $\varphi:=\Fl^{X}_{t,0}$ and $\ell(\varphi):=\Fl^{\ell(X)}_{t,0}$. Then we use equation~\eqref{ellhom} to compute
\begin{align*}
\vpf^{k}(\ell(Y))\big(\gamma(t),[\ell(\varphi)\!\circ\!\sigma\!\circ\!\varphi^{-1}]^{k}_{\gamma(t)}\big)
&= \frac{\rm d}{{\rm d}s}\bigg|_{s=0}\!\!\big(\gamma(t),\big[\Fl^{\ell(Y)}_{s}\circ\ell(\varphi) \circ\sigma\circ\varphi^{-1}\circ\Fl^{Y}_{-s}\big]^{k}_{\gamma(t)}\big)
\\
&= \frac{\rm d}{{\rm d}s}\bigg|_{s=0}\!\!\big(\gamma(t),\!\big[\ell(\varphi) \!\circ\Fl^{\ell(\varphi^{-1}_{*}(Y))}_{s}\! \circ\sigma\circ\Fl^{\varphi^{-1}_{*}(Y)}_{-s} \circ\varphi^{-1}\big]^{k}_{\gamma(t)}\big)
\\
&= 0
\end{align*}
by the $\FF$-invariance of $\sigma$ and since $\varphi^{-1}_{*}(Y)\in\FF$. Therefore equation~\eqref{ivpsol} does indeed define a curve in $\Gamma_{k}(\pi_{B})^{\FF}$. That equation~\eqref{ivpsol} defines a solution to the initial value problem of~equation~\eqref{ivp} follows by definition of the prolongation.

Uniqueness for $k$ finite follows from uniqueness of the flow in a manifold (since each $\Gamma_{k}(\pi_{B})^{\FF}$ is a subset of the jet manifold $\Gamma_{k}(\pi_{B})$), while for $k=\infty$ any solution to equation~\eqref{ivp} is the projective limit of the solutions for finite $k$, so uniqueness in this case follows from uniqueness for each finite $k$. Since $\Gamma_{\g}(\pi_{B})$ is neither a manifold nor a projective limit of manifolds, the argument is more subtle in the $k=\g$ case. Suppose that $\rho\colon [0,d]\rightarrow\Gamma_{\g}(\pi_{B})^{\FF}$ is some other solution to the initial value problem given in equation~\eqref{ivp}; thus in particular $\rho(0) = \big(x,[\sigma]^{\g}_{x}\big)$. Let us consider $\epsilon>0$ sufficiently small that for all $t\in[0,\epsilon)$ we can write $\rho(t) = \big(\gamma(t),[\sigma_{t}]^{\g}_{\gamma(t)}\big)$ for some $\FF$-invariant family $t\mapsto\sigma_{t}$ in $\Gamma_{\loc}(\pi_{B})$. By definition of the diffeology on $\Gamma_{\loc}(\pi_{B})$ we may assume that $\epsilon$ is sufficiently small that there is some open neighbourhood $\OO$ of $\gamma(0)$ containing $\gamma([0,\epsilon))$ and on which $\sigma_{t}$ is defined for all $t\in[0,\epsilon)$. Let us also assume that there are fibre bundle coordinates $\big(x^{1},\dots,x^{n};b^{1},\dots,b^{m}\big)$ on $B$ about $\sigma(x)$, the projection of whose domain to~$M$ contains $\OO$. By hypothesis we have\vspace{-1ex}
\begin{gather*}
\frac{\rm d}{{\rm d}s}\bigg|_{t}[\sigma_{s}]^{\g}_{\gamma(s)} = \frac{\rm d}{{\rm d}s}\bigg|_{t}\big[\Fl^{\ell(X)}_{s,0}\circ\sigma\circ\Fl^{X}_{0,s}\big]_{\gamma(s)}
\end{gather*}
for all $t\in[0,\epsilon)$, which by Proposition~\ref{tangentgerm} implies the coordinate expression\vspace{-1ex}
\begin{gather*}
\frac{\rm d}{{\rm d}s}\bigg|_{t}\big(\sigma^{i}_{s}(x^{1},\dots,x^{n})- \big(\Fl^{\ell(X)}_{s,0}\circ\sigma\circ\Fl^{X}_{0,s}\big)^{i}\big(x^{1},\dots,x^{n}\big)\big) = 0.
\end{gather*}
for $i=1,\dots,m$ and for $\big(x^{1},\dots,x^{n}\big)\in\OO$. This being true for all $t\in[0,\epsilon)$, it follows that the difference $\sigma^{i}_{t}-\big(\Fl^{\ell(X)}_{t,0}\circ\sigma\circ\Fl^{X}_{0,t}\big)^{i}$ is constant in $t$ on $\OO$ for all $i$, and since $\sigma^{i}_{0} = \sigma^{i}$ on $\OO$ we obtain $[\sigma_{t}]^{\g}_{\gamma(t)} = \big[\Fl^{\ell(X)}_{t,0}\circ\sigma\circ\Fl^{X}_{0,t}\big]^{\g}_{\gamma(t)}$ for all $t\in[0,\epsilon)$. Uniqueness on $[0,d]$ now follows from the compactness of $[0,d]$.


It remains only to show smoothness of $L\big(\pi_{B}^{k,\FF}\big)$. Let $\varphi_{\PP}:=\big(\tilde{\gamma},[\tilde{X}]^{\g},\tilde{d}\big)\colon U\rightarrow\PP(\FF)$ and $\varphi_{\Gamma}:=\big(\tilde{x},[\tilde{\sigma}]^{k}_{\tilde{x}}\big)\colon V\rightarrow\Gamma_{k}(\pi_{B})^{\FF}$ be plots. We~need to show that the map from $W:=U\times_{s\circ\varphi_{\PP},\pi_{B}^{k,\FF}\circ\varphi_{\Gamma}}V\times\RB_{\geq0}$ to $\Gamma_{k}(\pi_{B})^{\FF}$ defined by
\begin{gather*}
(u,v,t)\mapsto \big(\tilde{\gamma}(u)(t),\big[\Fl^{\ell(\tilde{X}(u))}_{t,0}\circ \tilde{\sigma}(v)\circ\Fl^{\tilde{X}(u)}_{0,t}\big]^{k}_{\tilde{\gamma}(u)(t)}\big)
\end{gather*}
is smooth. The~map $\tilde{\gamma}$ is already a plot of $\PP(M)$ by definition. Recalling that $(\Gamma_{\pi_{B}})_{\loc}$ denotes the space of all locally defined sections of $\pi_{B}$ equipped with the diffeology of Proposition~\ref{functprop}, it~suffices to show that the map
\begin{gather*}
(u,v,t)\mapsto\kappa(u,v,t):=\big(\Fl^{\ell(\tilde{X}(u))}_{t,0}\circ\tilde{\sigma}(v)\circ\Fl^{\tilde{X}(u)}_{0,t}\big)\in(\Gamma_{\pi_{B}})_{\loc}
\end{gather*}
is smooth. That is, fixing $(u_{0},v_{0},t_{0})\in W$ and $x_{0}\in\dom(\kappa(u_{0},v_{0},t_{0}))$, we must find an open neighbourhood $\WW$ of $(u_{0},v_{0},t_{0})$ in $W$ and an open neighbourhood $\OO$ of $x_{0}$ in $M$ such that $\OO\subset\dom(\kappa(u,v,t))$ for all $(u,v,t)\in\WW$ and for which the map
{\samepage\begin{gather*}
\WW\times\OO\ni(u,v,t,x)\mapsto \kappa(u,v,t)(x)\in B
\end{gather*}
is smooth in the usual sense. Setting $d_{0}:=d(u_{0})$, we have the following.

}

\begin{enumerate}\itemsep=0pt
\item Using Lemma~\ref{technicallemma} together with the smoothness of $\ell$ as a map $\Gamma_{\loc}(\FF)\rightarrow(\pi_{B})_{*}(\XF_{B})_{\loc}$, we can find open neighbourhoods $\OO^{B}\ni\tilde{\sigma}(v_{0})\big(\Fl^{\tilde{X}(u_{0})}_{0,t_{0}}(x_{0})\big)$ in $B$, $\UU_{1}\ni u_{0}$ in $U$ and $I_{1}\supset[0,t_{0}]$ in $\RB_{\geq0}$ for which $\OO^{B}\subset\dom\big(\Fl^{\ell(\tilde{X}(u))}_{t,0}\big)$ for all $(u,t)\in\UU_{1}\times I_{1}$, and such that $I_{1}\times\UU_{1}\times \OO^{B}\ni(t,u,b)\mapsto\Fl^{\ell(\tilde{X}(u))}_{t,0}(b)\in B$ is smooth.
\item By definition of the diffeology on $(\Gamma_{\pi_{B}})_{\loc}$, we can find open neighbourhoods $\OO^{M}\ni\Fl^{\tilde{X}(u_{0})}_{0,t_{0}}(x_{0})$ in $M$ and $\VV\ni v_{0}$ in $V$ such that $\OO^{M}\subset\dom(\tilde{\sigma}(v))$ and $\tilde{\sigma}(v)(\OO^{M})\subset\OO^{B}$ for all $v\in\VV$, and for which $\VV\times\OO^{M}\ni(v,x)\mapsto\tilde{\sigma}(v)(x)\in B$ is smooth.
\item Again by Lemma~\ref{technicallemma}, we can find open neighbourhoods $\UU_{2}\ni u_{0}$, $I_{2}\supset[0,t_{0}]$ in $\RB_{\geq0}$ and $\OO\ni x_{0}$ in $M$ such that $\OO\subset\dom\big(\Fl^{\tilde{X}(u)}_{0,t}\big)$ and $\Fl^{\tilde{X}(u)}_{0,t}(\OO)\subset\OO^{M}$ for all $u\in\UU_{2}$ and $t\in I_{2}$, and such that $I_{2}\times\UU_{2}\times \OO\ni(t,u,x)\mapsto\Fl^{\tilde{X}(u)}_{0,t}(x)\in M$ is smooth.
\end{enumerate}
Finally, therefore, setting $\UU:=\UU_{1}\cap\UU_{2}$, $I:=I_{1}\cap I_{2}$ and
\begin{gather*}
\WW:=\UU\times_{s\circ\varphi_{\PP},\pi_{B}^{k,\FF}\circ\varphi_{\Gamma}}\VV\times I\subset W,
\end{gather*}
we have $\OO\subset\dom(\kappa(u,v,t))$ for all $(u,v,t)\in\WW$ and $\WW\times\OO\ni(u,v,t,x)\mapsto\kappa(u,v,t)(x)\in B$ is smooth.
\end{proof}
\begin{thebibliography}{99}
\bibitem[1]{varbi}
Anderson I., The variational bicomplex, Utah State University, 1989.

\bibitem[2]{iakovos2}
Androulidakis I., Skandalis G., The holonomy groupoid of a singular foliation,
 \href{https://doi.org/10.1515/CRELLE.2009.001}{\textit{J.~Reine Angew. Math.}} \textbf{626} (2009), 1--37,
 \href{https://arxiv.org/abs/math.DG/0612370}{arXiv:math.DG/0612370}.

\bibitem[7]{aac1}
Arias~Abad C., Crainic M., The {W}eil algebra and the {V}an {E}st isomorphism,
 \href{https://doi.org/10.5802/aif.2633}{\textit{Ann. Inst. Fourier (Grenoble)}} \textbf{61} (2011), 927--970,
 \href{https://arxiv.org/abs/0901.0322}{arXiv:0901.0322}.

\bibitem[14]{bryantgriffiths1}
Bryant R.L., Griffiths P.A., Characteristic cohomology of differential systems.
 {I}.~{G}eneral theory, \href{https://doi.org/10.2307/2152923}{\textit{J.~Amer. Math. Soc.}} \textbf{8} (1995),
 507--596.

\bibitem[16]{cw}
Christensen J.D., Wu E., Tangent spaces and tangent bundles for diffeological
 spaces, \textit{Cah. Topol. G\'eom. Diff\'er. Cat\'eg.} \textbf{57} (2016),
 3--50, \href{https://arxiv.org/abs/1411.5425}{arXiv:1411.5425}.

\bibitem[31]{GarZam}
Garmendia A., Zambon M., Hausdorff {M}orita equivalence of singular foliations,
 \href{https://doi.org/10.1007/s10455-018-9620-6}{\textit{Ann. Global Anal. Geom.}} \textbf{55} (2019), 99--132,
 \href{https://arxiv.org/abs/1803.00896}{arXiv:1803.00896}.

\bibitem[36]{hartshorne}
Hartshorne R., Algebraic geometry, \textit{Graduate Texts in Mathematics},
 Vol.~52, \href{https://doi.org/10.1007/978-1-4757-3849-0}{Springer-Verlag}, New York~-- Heidelberg, 1977.

\bibitem[37]{hector}
Hector G., G\'eom\'etrie et topologie des espaces diff\'eologiques, in Analysis
 and Geometry in Foliated Manifolds ({S}antiago de {C}ompostela, 1994), World
 Sci. Publ., River Edge, NJ, 1995, 55--80.

\bibitem[40]{diffeology}
Iglesias-Zemmour P., Diffeology, \textit{Mathematical Surveys and Monographs},
 Vol.~185, \href{https://doi.org/10.1090/surv/185}{Amer. Math. Soc.}, Providence, RI, 2013.

\bibitem[41]{folbund}
Kamber F.W., Tondeur P., Foliated bundles and characteristic classes,
 \textit{Lecture Notes in Math.}, Vol.~493, \href{https://doi.org/10.1007/BFb0081558}{Springer-Verlag}, Berlin~-- New
 York, 1975.

\bibitem[42]{natopdiffgeom}
Kol\'a\v{r} I., Michor P.W., Slov\'ak J., Natural operations in differential
 geometry, \href{https://doi.org/10.1007/978-3-662-02950-3}{Springer-Verlag}, Berlin, 1993.

\bibitem[44]{leesmth}
Lee J.M., Introduction to smooth manifolds, 2nd ed., \textit{Graduate Texts in
 Mathematics}, Vol.~218, \href{https://doi.org/10.1007/978-0-387-21752-9}{Springer}, New York, 2013.

\bibitem[47]{mac3}
MacDonald L.E., Hierarchies of holonomy groupoids for foliated bundles,
 \href{https://arxiv.org/abs/2004.13929}{arXiv:2004.13929}.

\bibitem[52]{olver}
Olver P.J., Applications of {L}ie groups to differential equations,
 \textit{Graduate Texts in Mathematics}, Vol.~107, \href{https://doi.org/10.1007/978-1-4684-0274-2}{Springer-Verlag}, New York,
 1986.

\bibitem[56]{saunders}
Saunders D.J., The geometry of jet bundles, \textit{London Mathematical Society
 Lecture Note Series}, Vol.~142, \href{https://doi.org/10.1017/CBO9780511526411}{Cambridge University Press}, Cambridge, 1989.

\bibitem[57]{ptfunctor}
Schreiber U., Waldorf K., Parallel transport and functors, \textit{J.~Homotopy
 Relat. Struct.} \textbf{4} (2009), 187--244, \href{https://arxiv.org/abs/0705.0452}{arXiv:0705.0452}.

\bibitem[58]{stefan}
Stefan P., Accessible sets, orbits, and foliations with singularities,
 \href{https://doi.org/10.1112/plms/s3-29.4.699}{\textit{Proc. London Math. Soc.}} \textbf{29} (1974), 699--713.

\bibitem[59]{sussmann}
Sussmann H.J., Orbits of families of vector fields and integrability of
 distributions, \href{https://doi.org/10.2307/1996660}{\textit{Trans. Amer. Math. Soc.}} \textbf{180} (1973),
 171--188.

\bibitem[60]{VilGar1}
Villatoro~J. G.A., Integration of singular foliations via paths,
 \href{https://arxiv.org/abs/1912.02148}{arXiv:1912.02148}.

\end{thebibliography}
\end{document}
